% Generated by tools/build.py -- do not edit; edit data/computer-networks/*.txt
% No preamble here: this file is \input by main.tex.
\chapter{Computer Networks}

\section{Data Communication (signals, channel capacity, encoding, multiplexing, media, error handling)}

\begin{question}{H}
A noiseless channel of bandwidth 3000 Hz transmits a signal with two signal levels. The maximum bit rate of the channel is
\optsshort{1500 bps}{3000 bps}{6000 bps}{12000 bps}
\end{question}

\begin{question}{H}
A channel has a bandwidth of 4 kHz and a signal-to-noise ratio of 15. The maximum achievable data rate is
\optsshort{8 kbps}{12 kbps}{16 kbps}{60 kbps}
\end{question}

\begin{question}{H}
What is the approximate capacity of a telephone line with a bandwidth of 3 kHz and a signal-to-noise ratio of 30 dB?
\optsshort{90 kbps}{10 kbps}{30 kbps}{3 kbps}
\end{question}

\begin{question}{M}
A signal-to-noise ratio of 20 dB corresponds to a power ratio of
\optsshort{20}{10}{100}{200}
\end{question}

\begin{question}{H}
A modem uses 16-QAM and transmits at a bit rate of 4800 bps. The baud rate is
\optsshort{19200 baud}{300 baud}{4800 baud}{1200 baud}
\end{question}

\begin{question}{H}
What is the baud rate of standard 10 Mbps Ethernet that uses Manchester encoding?
\optsshort{40 Mbaud}{10 Mbaud}{20 Mbaud}{5 Mbaud}
\end{question}

\begin{question}{M}
The propagation delay of a 2000 km link in which the signal travels at 2 $\times$ 10\textsuperscript{8} m/s is
\optsshort{100 ms}{1 ms}{10 ms}{4 ms}
\end{question}

\begin{question}{M}
The transmission time of a 1000-byte frame over a 1 Mbps link is
\optsshort{1 s}{8 ms}{125 $\mu$s}{1 ms}
\end{question}

\begin{question}{H}
A block code has a minimum Hamming distance of 5. It can
\optslong{detect up to 5 errors and correct up to 2 errors}{detect up to 4 errors and correct up to 2 errors}{detect up to 4 errors and correct up to 4 errors}{detect up to 2 errors and correct up to 1 error}
\end{question}

\begin{question}{H}
How many redundant bits are needed in a Hamming code to correct a single-bit error in a 7-bit data word?
\optsshort{2}{4}{5}{3}
\end{question}

\begin{question}{H}
A bit stream 1101011011 is transmitted using the CRC method with generator polynomial x\textsuperscript{4} + x + 1. The CRC bits appended to the message are
\optsshort{0111}{1001}{1011}{1110}
\end{question}

\begin{question}{H}
The message 100100 is to be transmitted using CRC with generator 1101. The transmitted frame is
\optsshort{100100000}{100100001}{100100110}{100100011}
\end{question}

\begin{question}{M}
Using even parity, the parity bit appended to the data unit 1011001 is
\optsshort{0}{either 0 or 1}{cannot be determined}{1}
\end{question}

\begin{question}{H}
Five channels, each with a 100 kHz bandwidth, are to be frequency-division multiplexed with guard bands of 10 kHz between adjacent channels. The minimum bandwidth of the link is
\optsshort{550 kHz}{600 kHz}{540 kHz}{500 kHz}
\end{question}

\begin{question}{H}
In bit-oriented framing with flag 01111110, the data 0111111011111 is transmitted after bit stuffing as
\optsshort{011111101111101}{011111011011111}{01111101011111}{011111010111110}
\end{question}

\begin{question}{H}
A voice channel is band-limited to 4 kHz and digitised using PCM with 256 quantisation levels. The bit rate produced is
\optsshort{56 kbps}{32 kbps}{128 kbps}{64 kbps}
\end{question}

\begin{question}{M}
A signal loses half its power while travelling through a medium. The attenuation is approximately
\optsshort{$-$10 dB}{$-$1 dB}{$-$3 dB}{$-$2 dB}
\end{question}

\begin{question}{M}
The T1 carrier multiplexes 24 voice channels and has a data rate of
\optsshort{2.048 Mbps}{1.544 Mbps}{6.312 Mbps}{64 kbps}
\end{question}

\begin{question}{M}
Light is guided through an optical fibre by the principle of
\optsshort{total internal reflection}{refraction at the core}{dispersion}{diffraction}
\end{question}

\begin{question}{H}
Match List I with List II.\\
List I (Line coding scheme): A. NRZ-I  B. Manchester  C. AMI  D. Differential Manchester\\
List II (Property): I. Bipolar; alternate 1s take opposite polarity  II. Transition at the middle of every bit; the direction of the transition encodes the bit  III. Inversion at the start of a bit denotes 1  IV. Mid-bit transition always present; presence or absence of a transition at the start of the bit encodes the bit
\optsshort{A-III, B-IV, C-I, D-II}{A-II, B-III, C-I, D-IV}{A-III, B-II, C-I, D-IV}{A-I, B-II, C-III, D-IV}
\end{question}

\begin{question}{H}
Given below are two statements: one is labelled as Assertion A and the other is labelled as Reason R.\\
Assertion A: Manchester encoding provides self-synchronisation between sender and receiver.\\
Reason R: In Manchester encoding there is a transition in the middle of every bit interval.
\optslong{Both A and R are true and R is the correct explanation of A}{Both A and R are true but R is NOT the correct explanation of A}{A is true but R is false}{A is false but R is true}
\end{question}

\begin{question}{H}
Given below are two statements.\\
Statement I: The Nyquist bit rate formula gives the theoretical maximum bit rate of a noiseless channel.\\
Statement II: According to the Shannon capacity formula, the capacity of a noisy channel can be increased without limit by increasing the number of signal levels.
\optslong{Both Statement I and Statement II are true}{Both Statement I and Statement II are false}{Statement I is true but Statement II is false}{Statement I is false but Statement II is true}
\end{question}

\begin{question}{M}
Match List I with List II.\\
List I (Device/system): A. Keyboard to CPU  B. Walkie-talkie  C. Telephone network\\
List II (Transmission mode): I. Full-duplex  II. Simplex  III. Half-duplex
\optsshort{A-I, B-III, C-II}{A-II, B-III, C-I}{A-III, B-II, C-I}{A-II, B-I, C-III}
\end{question}

\begin{question}{M}
Which of the following is NOT one of the five components of a data communication system?
\optsshort{Transmission medium}{Message}{Protocol}{Topology}
\end{question}

\begin{question}{M}
How many bits are carried by each signal element in QPSK?
\optsshort{3}{1}{4}{2}
\end{question}

\begin{question}{M}
A digital signal has eight levels. How many bits are needed per level?
\optsshort{3}{8}{2}{4}
\end{question}

\begin{question}{M}
A periodic signal is decomposed into five sine waves with frequencies of 100, 300, 500, 700 and 900 Hz. Its bandwidth is
\optsshort{500 Hz}{900 Hz}{800 Hz}{2500 Hz}
\end{question}

\begin{question}{L}
A sine wave has a period of 1 ms. Its frequency is
\optsshort{1 Hz}{1 MHz}{100 Hz}{1 kHz}
\end{question}

\begin{question}{H}
A channel has a bandwidth of 1 MHz and a signal-to-noise ratio of 63. The appropriate upper limit on the bit rate is
\optsshort{1 Mbps}{6 Mbps}{4 Mbps}{63 Mbps}
\end{question}

\begin{question}{M}
In which multiplexing technique are time slots allocated dynamically, on demand, so that no slot is wasted on an idle source?
\optsshort{Statistical TDM}{Synchronous TDM}{FDM}{WDM}
\end{question}

\begin{question}{M}
A two-dimensional parity check can
\optslong{detect only single-bit errors}{detect and correct all 2-bit errors}{detect all errors of up to 3 bits and correct single-bit errors}{correct all burst errors}
\end{question}

\begin{question}{M}
In asynchronous serial transmission each 8-bit character is sent with one start bit and one stop bit. The transmission efficiency is
\optsshort{70\%}{100\%}{88.9\%}{80\%}
\end{question}

\begin{question}{M}
In a 7-bit Hamming code (4 data bits), the parity bits are placed at positions
\optsshort{1, 2 and 3}{1, 3 and 5}{5, 6 and 7}{1, 2 and 4}
\end{question}

\begin{question}{H}
According to the Nyquist sampling theorem, an analog signal whose highest frequency is f must be sampled at a rate of at least
\optsshort{f / 2}{f}{2f}{4f}
\end{question}

\begin{question}{M}
The wires in a twisted-pair cable are twisted together in order to
\optslong{increase the tensile strength of the cable}{reduce crosstalk and electromagnetic interference}{increase the data rate by doubling the conductors}{reduce the cost of the cable}
\end{question}

\begin{question}{M}
A geostationary satellite orbits at an altitude of approximately
\optsshort{20,200 km}{2,000 km}{72,000 km}{35,786 km}
\end{question}

\begin{question}{M}
Given below are two statements.\\
Statement I: In frequency hopping spread spectrum (FHSS) the carrier frequency hops according to a pseudo-random sequence known to both sender and receiver.\\
Statement II: In direct sequence spread spectrum (DSSS) each data bit is replaced by a sequence of chips.
\optslong{Both Statement I and Statement II are true}{Both Statement I and Statement II are false}{Statement I is true but Statement II is false}{Statement I is false but Statement II is true}
\end{question}

\begin{question}{M}
Given below are two statements.\\
Statement I: Baseband transmission sends a digital signal over the channel without changing it to an analog signal.\\
Statement II: Broadband transmission requires a low-pass channel starting at zero frequency.
\optslong{Both Statement I and Statement II are true}{Both Statement I and Statement II are false}{Statement I is true but Statement II is false}{Statement I is false but Statement II is true}
\end{question}

\begin{question}{M}
Wavelength-division multiplexing is conceptually the same as
\optslong{TDM applied to optical signals}{FDM applied to optical signals}{statistical TDM}{CDM applied to optical signals}
\end{question}

\begin{question}{H}
A CRC generator polynomial of degree r detects
\optslong{only burst errors of length exactly r}{all burst errors of length less than or equal to r}{all burst errors of any length}{only single-bit errors}
\end{question}

\begin{question}{M}
Arrange the following guided media in increasing order of bandwidth.\\
A. Coaxial cable  B. Optical fibre  C. Twisted-pair cable
\optsshort{C, B, A}{C, A, B}{A, C, B}{B, A, C}
\end{question}

\begin{question}{M}
Which of the following digital-to-analog modulation techniques is the most susceptible to noise?
\optsshort{FSK}{PSK}{ASK}{QPSK}
\end{question}

\begin{question}{H}
Quadrature amplitude modulation (QAM) is a combination of
\optsshort{ASK and FSK}{PSK and PCM}{ASK and PSK}{FSK and PSK}
\end{question}

\begin{question}{M}
How many bits per sample are transmitted in delta modulation?
\optslong{8}{2}{1}{It depends on the number of quantisation levels}
\end{question}

\begin{question}{M}
Block coding scheme 4B/5B is used before NRZ-I in 100Base-TX mainly to
\optslong{double the data rate}{ensure enough transitions for synchronisation}{correct single-bit errors}{compress the data}
\end{question}

\begin{question}{M}
B8ZS and HDB3 are techniques used for
\optslong{scrambling long runs of zeros in AMI}{multiplexing}{analog-to-digital conversion}{error correction}
\end{question}

\begin{question}{H}
The 8-bit one's-complement checksum of the two data words 10101001 and 00111001 is
\optsshort{00011101}{11100011}{00011110}{11100010}
\end{question}

\begin{question}{M}
Given below are two statements.\\
Statement I: Throughput of a link can never be greater than its bandwidth in bits per second.\\
Statement II: Latency of a network is the sum of propagation time, transmission time, queuing time and processing delay.
\optslong{Both Statement I and Statement II are true}{Both Statement I and Statement II are false}{Statement I is true but Statement II is false}{Statement I is false but Statement II is true}
\end{question}

\begin{question}{M}
Match List I with List II.\\
List I (Impairment): A. Attenuation  B. Distortion  C. Crosstalk  D. Thermal noise\\
List II (Description): I. Signal components travel at different speeds and arrive out of phase  II. Random motion of electrons in the conductor  III. Loss of signal energy with distance  IV. Coupling of a signal from one wire onto another
\optsshort{A-I, B-III, C-IV, D-II}{A-II, B-I, C-IV, D-III}{A-III, B-I, C-IV, D-II}{A-III, B-IV, C-I, D-II}
\end{question}

\begin{question}{M}
A signal is transmitted at 1000 baud using 8-PSK. The bit rate is
\optsshort{1000 bps}{125 bps}{3000 bps}{8000 bps}
\end{question}

\begin{question}{H}
A noiseless 3 kHz channel has to carry 24 kbps. The minimum number of signal levels required is
\optsshort{256}{4}{16}{8}
\end{question}

\begin{question}{M}
Four sources, each producing 1 kbps, are combined using synchronous TDM with one bit from each source per frame. The duration of a bit on the multiplexed link is
\optsshort{4 ms}{250 $\mu$s}{500 $\mu$s}{1 ms}
\end{question}

\begin{question}{M}
Given below are two statements.\\
Statement I: Single-mode fibre has a much smaller core diameter than multimode fibre.\\
Statement II: Single-mode fibre suffers from greater modal dispersion than multimode fibre.
\optslong{Both Statement I and Statement II are true}{Both Statement I and Statement II are false}{Statement I is true but Statement II is false}{Statement I is false but Statement II is true}
\end{question}

\begin{question}{M}
Match List I with List II.\\
List I (Unguided medium): A. Radio waves  B. Microwaves  C. Infrared\\
List II (Characteristic): I. Cannot penetrate walls; short-range communication  II. Omnidirectional propagation  III. Unidirectional, line-of-sight propagation
\optsshort{A-III, B-II, C-I}{A-I, B-III, C-II}{A-II, B-III, C-I}{A-II, B-I, C-III}
\end{question}

\begin{question}{M}
In forward error correction (FEC)
\optslong{the receiver corrects errors using redundant bits without retransmission}{the receiver asks the sender to retransmit corrupted frames}{errors are ignored}{only error detection is possible}
\end{question}

\begin{question}{H}
A 7-bit Hamming code word (even parity, bit positions numbered 1 to 7 from left to right, parity bits at positions 1, 2 and 4) is received as 1011011. The bit in error is at position
\optsshort{6}{5}{7}{3}
\end{question}

\begin{question}{L}
A device that converts digital data to an analog signal and back is called a
\optsshort{multiplexer}{codec}{modem}{repeater}
\end{question}

\begin{question}{M}
Which of the following multiplexing techniques is used for digital signals?
\optsshort{FDM}{WDM}{TDM}{Both FDM and WDM}
\end{question}

\begin{question}{M}
In CDMA, the chip sequences assigned to different stations are
\optslong{orthogonal to each other}{identical}{random and unknown to the receiver}{complements of each other}
\end{question}

\begin{question}{M}
The maximum data rate of a channel with bandwidth B and a signal-to-noise ratio of 0 is
\optsshort{2B}{infinite}{B}{0}
\end{question}

\begin{question}{M}
The E1 carrier has a data rate of
\optsshort{2.048 Mbps}{44.736 Mbps}{8.448 Mbps}{1.544 Mbps}
\end{question}

\begin{question}{L}
Which connector is used with unshielded twisted-pair cable in Ethernet LANs?
\optsshort{RJ-45}{SC}{RJ-11}{BNC}
\end{question}

\begin{question}{M}
Which of the following statements about a repeater and an amplifier is correct?
\optslong{A repeater amplifies the signal together with the noise}{An amplifier regenerates the original bit pattern}{A repeater regenerates the signal while an amplifier amplifies both signal and noise}{Both operate at the data link layer}
\end{question}

\begin{question}{H}
Given below are two statements: one is labelled as Assertion A and the other is labelled as Reason R.\\
Assertion A: A cyclic redundancy check can correct single-bit errors in a frame.\\
Reason R: CRC is based on modulo-2 binary division.
\optslong{Both A and R are true and R is the correct explanation of A}{Both A and R are true but R is NOT the correct explanation of A}{A is true but R is false}{A is false but R is true}
\end{question}

\begin{question}{M}
Which of the following is an example of a polar line coding scheme with no DC-component problem and self-synchronisation?
\optsshort{Manchester}{NRZ-I}{Unipolar NRZ}{NRZ-L}
\end{question}

\begin{question}{M}
The minimum bandwidth required for an ASK signal transmitting at 2000 bps (with d = 0, i.e. ideal filtering) in half-duplex mode is
\optsshort{4000 Hz}{1000 Hz}{2000 Hz}{8000 Hz}
\end{question}

\begin{question}{L}
In data communication, jitter refers to
\optsshort{variation in packet delay}{loss of packets}{the average delay of packets}{corruption of bits}
\end{question}

\begin{answers}
\ans{1}{3}{Nyquist: C = 2B log2 L = 2 $\times$ 3000 $\times$ 1 = 6000 bps.}
\ans{2}{3}{Shannon: C = B log2(1 + SNR) = 4000 $\times$ log2 16 = 16 kbps.}
\ans{3}{3}{30 dB = 1000; C = 3000 $\times$ log2(1001) $\approx$ 3000 $\times$ 9.97 $\approx$ 30 kbps.}
\ans{4}{3}{SNR(dB) = 10 log10(S/N), so S/N = 10\textsuperscript{2} = 100.}
\ans{5}{4}{16 = 2\textsuperscript{4}, so 4 bits per signal element; baud = 4800 / 4 = 1200.}
\ans{6}{3}{Manchester uses two signal elements per bit, so the baud rate is twice the bit rate.}
\ans{7}{3}{Tp = distance / speed = 2 $\times$ 10\textsuperscript{6} / 2 $\times$ 10\textsuperscript{8} = 10 ms.}
\ans{8}{2}{Tt = 8000 bits / 10\textsuperscript{6} bps = 8 ms.}
\ans{9}{2}{Detection of s errors needs dmin = s + 1; correction of t errors needs dmin = 2t + 1.}
\ans{10}{2}{Need 2\textsuperscript{r} $\geq$ m + r + 1; r = 4 gives 16 $\geq$ 12, while r = 3 gives 8 < 11.}
\ans{11}{4}{Dividing 11010110110000 by 10011 (modulo-2) leaves remainder 1110.}
\ans{12}{2}{100100000 divided by 1101 (modulo-2) gives remainder 001, which is appended to the data.}
\ans{13}{1}{The data already has four 1s (an even number), so the parity bit is 0.}
\ans{14}{3}{5 $\times$ 100 + 4 guard bands $\times$ 10 = 540 kHz.}
\ans{15}{4}{A 0 is stuffed after every run of five consecutive 1s: 011111 0 10 11111 0.}
\ans{16}{4}{Sampling at 2 $\times$ 4000 = 8000 samples/s with 8 bits per sample gives 64 kbps.}
\ans{17}{3}{10 log10(0.5) $\approx$ $-$3 dB.}
\ans{18}{2}{T1 = 24 $\times$ 8 bits + 1 framing bit = 193 bits per 125 $\mu$s = 1.544 Mbps. E1 is 2.048 Mbps.}
\ans{19}{1}{The core has a higher refractive index than the cladding, so light striking beyond the critical angle is totally reflected.}
\ans{20}{3}{NRZ-I inverts on 1; Manchester encodes by direction of the mid-bit transition; AMI is bipolar; differential Manchester uses the transition at the bit boundary.}
\ans{21}{1}{The guaranteed mid-bit transition lets the receiver recover the clock, which is exactly why the code is self-clocking.}
\ans{22}{3}{Shannon capacity depends only on bandwidth and SNR; the number of levels does not appear in it.}
\ans{23}{2}{Keyboard input is one-way (simplex), a walkie-talkie alternates direction (half-duplex), telephony is simultaneous both ways (full-duplex).}
\ans{24}{4}{The five components are message, sender, receiver, transmission medium and protocol.}
\ans{25}{4}{QPSK uses four phases, and log2 4 = 2 bits per signal element.}
\ans{26}{1}{log2 8 = 3 bits per level.}
\ans{27}{3}{Bandwidth = highest $-$ lowest frequency = 900 $-$ 100 = 800 Hz.}
\ans{28}{4}{f = 1/T = 1 / 10\textsuperscript{-3} = 1000 Hz.}
\ans{29}{2}{C = 10\textsuperscript{6} $\times$ log2(64) = 6 Mbps.}
\ans{30}{1}{Statistical (asynchronous) TDM allocates slots only to sources that have data, so each slot carries an address.}
\ans{31}{3}{Row and column parity locates a single-bit error (so it can be corrected) and detects any 1-, 2- or 3-bit error.}
\ans{32}{4}{8 data bits out of 10 transmitted bits = 80\%.}
\ans{33}{4}{Parity bits occupy the positions that are powers of two: 1, 2, 4.}
\ans{34}{3}{Sampling at least twice the highest frequency allows exact reconstruction.}
\ans{35}{2}{Twisting makes both wires pick up nearly equal noise, which cancels at the receiver.}
\ans{36}{4}{GEO satellites orbit at about 35,786 km above the equator with a 24-hour period.}
\ans{37}{1}{Both are the defining properties of the two spread-spectrum techniques.}
\ans{38}{3}{Broadband transmission uses modulation and therefore a band-pass channel; baseband needs the low-pass channel.}
\ans{39}{2}{WDM combines light sources of different wavelengths (frequencies) on one fibre, i.e. FDM at optical frequencies.}
\ans{40}{2}{All bursts of length $\leq$ r are detected; longer bursts are detected with high probability.}
\ans{41}{2}{Twisted pair < coaxial cable < optical fibre.}
\ans{42}{3}{Noise mostly affects amplitude, and ASK carries its information in the amplitude.}
\ans{43}{3}{QAM varies both the amplitude and the phase of the carrier.}
\ans{44}{3}{Delta modulation sends a single bit indicating whether the signal rose or fell since the previous sample.}
\ans{45}{2}{4B/5B guarantees no long run of zeros, giving NRZ-I enough transitions; it adds 25\% overhead.}
\ans{46}{1}{Both substitute a run of zeros (8 and 4 respectively) with a pattern containing violations to keep synchronisation.}
\ans{47}{1}{The sum is 11100010; its one's complement is 00011101.}
\ans{48}{1}{Bandwidth is the potential rate, throughput is the achieved rate; latency is the sum of the four delay components.}
\ans{49}{3}{Attenuation is energy loss, distortion is change of shape, crosstalk is coupling between wires, thermal noise is from electron motion.}
\ans{50}{3}{8-PSK carries log2 8 = 3 bits per baud, so 3000 bps.}
\ans{51}{3}{24000 = 2 $\times$ 3000 $\times$ log2 L gives log2 L = 4, so L = 16.}
\ans{52}{2}{The link carries 4 kbps, so each bit lasts 1/4000 s = 250 $\mu$s.}
\ans{53}{3}{Single-mode fibre carries one mode, so it has essentially no modal dispersion and supports longer distances.}
\ans{54}{3}{Radio waves are omnidirectional, microwaves need line of sight, infrared is confined to a room.}
\ans{55}{1}{FEC adds enough redundancy for the receiver to correct errors itself; ARQ uses retransmission.}
\ans{56}{3}{c1 = b1$\oplus$b3$\oplus$b5$\oplus$b7 = 1, c2 = b2$\oplus$b3$\oplus$b6$\oplus$b7 = 1, c4 = b4$\oplus$b5$\oplus$b6$\oplus$b7 = 1; syndrome 111 = 7.}
\ans{57}{3}{A modem (modulator-demodulator) does digital-to-analog conversion; a codec converts an analog signal to digital form.}
\ans{58}{3}{FDM and WDM are analog techniques; TDM is a digital multiplexing technique.}
\ans{59}{1}{The inner product of two different chip sequences is 0, which lets the receiver separate the stations.}
\ans{60}{4}{C = B log2(1 + 0) = 0; nothing can be received when the signal is completely lost in noise.}
\ans{61}{1}{E1 has 32 slots of 64 kbps each (30 for voice) = 2.048 Mbps.}
\ans{62}{1}{UTP Ethernet uses the 8-pin RJ-45 connector; BNC is used with coaxial cable and SC with fibre.}
\ans{63}{3}{A repeater is a regenerator at the physical layer; an amplifier cannot distinguish signal from noise.}
\ans{64}{4}{CRC is an error-detecting code only; it is indeed computed by modulo-2 division.}
\ans{65}{1}{Manchester has a transition in every bit, so there is no baseline wandering or DC component.}
\ans{66}{3}{For binary ASK, B = (1 + d) $\times$ S, with S = 2000 baud, so B = 2000 Hz.}
\ans{67}{1}{Jitter is the variation in the arrival delay of packets belonging to the same stream.}
\end{answers}

\section{Network Models, Topologies, Data Link Layer, MAC and LANs}

\begin{question}{H}
Match List I with List II.\\
List I (OSI layer): A. Physical  B. Data link  C. Network  D. Transport\\
List II (Protocol data unit): I. Segment  II. Bits  III. Packet  IV. Frame
\optsshort{A-II, B-IV, C-I, D-III}{A-II, B-IV, C-III, D-I}{A-II, B-III, C-IV, D-I}{A-IV, B-II, C-III, D-I}
\end{question}

\begin{question}{H}
Match List I with List II.\\
List I (Device): A. Repeater  B. Bridge  C. Router  D. Gateway\\
List II (Highest OSI layer of operation): I. Network layer  II. All seven layers  III. Physical layer  IV. Data link layer
\optsshort{A-III, B-I, C-IV, D-II}{A-III, B-IV, C-I, D-II}{A-IV, B-III, C-I, D-II}{A-III, B-IV, C-II, D-I}
\end{question}

\begin{question}{H}
How many point-to-point links are required to connect 10 devices in a fully connected mesh topology?
\optsshort{20}{90}{45}{10}
\end{question}

\begin{question}{M}
In a fully connected mesh network of n devices, the number of I/O ports required on each device is
\optsshort{2n}{n $-$ 1}{n}{n(n $-$ 1)/2}
\end{question}

\begin{question}{H}
Which layer of the OSI model is responsible for dialog control and synchronisation (check-pointing)?
\optsshort{Transport layer}{Data link layer}{Presentation layer}{Session layer}
\end{question}

\begin{question}{H}
Data compression, encryption and translation between data formats are functions of the
\optsshort{presentation layer}{application layer}{transport layer}{session layer}
\end{question}

\begin{question}{H}
Which of the following is the correct order of encapsulation as data moves down the layers of the sending host?
\optslong{Data, segment, packet, frame, bits}{Data, packet, segment, frame, bits}{Bits, frame, packet, segment, data}{Data, frame, packet, segment, bits}
\end{question}

\begin{question}{M}
Which OSI layers are combined into the application layer of the TCP/IP protocol suite?
\optslong{Application and presentation only}{Application, presentation and session}{Session and transport}{Presentation, session and transport}
\end{question}

\begin{question}{H}
Match List I with List II.\\
List I (Address): A. Port address  B. IPv4 address  C. MAC address  D. IPv6 address\\
List II (Size): I. 48 bits  II. 128 bits  III. 16 bits  IV. 32 bits
\optsshort{A-III, B-IV, C-II, D-I}{A-III, B-IV, C-I, D-II}{A-III, B-I, C-IV, D-II}{A-IV, B-III, C-I, D-II}
\end{question}

\begin{question}{H}
In a stop-and-wait protocol, the frame transmission time is 1 ms and the one-way propagation delay is 2 ms. Ignoring ACK size and processing time, the link utilisation is
\optsshort{20\%}{33.3\%}{25\%}{50\%}
\end{question}

\begin{question}{H}
A stop-and-wait protocol sends 1000-bit frames over a 1 Mbps link whose one-way propagation delay is 4.5 ms. Ignoring acknowledgement transmission time, the throughput is
\optsshort{200 kbps}{100 kbps}{500 kbps}{1 Mbps}
\end{question}

\begin{question}{H}
With a 3-bit sequence number field, the maximum sender window sizes in Go-Back-N and Selective Repeat are respectively
\optsshort{7 and 7}{8 and 8}{7 and 4}{8 and 4}
\end{question}

\begin{question}{H}
In the Selective Repeat ARQ protocol the sender window size is 32. The minimum number of bits needed for the sequence number field is
\optsshort{5}{4}{6}{7}
\end{question}

\begin{question}{H}
The receiver window size in the Go-Back-N ARQ protocol is
\optsshort{2\textsuperscript{m} $-$ 1}{2\textsuperscript{m}}{1}{2\textsuperscript{m$-$1}}
\end{question}

\begin{question}{M}
In a sliding window protocol, the optimal window size for maximum utilisation, where a = Tp/Tt, is
\optsshort{2a}{1 + a}{1 + 2a}{a}
\end{question}

\begin{question}{H}
The maximum throughput of pure ALOHA and slotted ALOHA are respectively
\optsshort{50\% and 100\%}{18.4\% and 50\%}{36.8\% and 18.4\%}{18.4\% and 36.8\%}
\end{question}

\begin{question}{H}
A pure ALOHA network transmits 200-bit frames on a shared 200 kbps channel. If the stations together produce 1000 frames per second, the throughput is approximately
\optsshort{184 frames per second}{500 frames per second}{368 frames per second}{135 frames per second}
\end{question}

\begin{question}{M}
The vulnerable time in pure ALOHA, where Tfr is the frame transmission time, is
\optsshort{Tfr / 2}{2 $\times$ Tfr}{Tfr}{4 $\times$ Tfr}
\end{question}

\begin{question}{H}
A CSMA/CD network runs at 10 Mbps over a cable of length 2500 m with a signal speed of 2 $\times$ 10\textsuperscript{8} m/s. The minimum frame size is
\optsshort{250 bits}{125 bits}{512 bits}{500 bits}
\end{question}

\begin{question}{H}
The minimum and maximum frame lengths (excluding preamble and SFD) in standard IEEE 802.3 Ethernet are
\optsshort{72 bytes and 1526 bytes}{64 bytes and 1518 bytes}{64 bytes and 1500 bytes}{46 bytes and 1500 bytes}
\end{question}

\begin{question}{H}
In the binary exponential backoff algorithm of CSMA/CD, after the 3rd collision a station waits a random number of slot times chosen from
\optsshort{0 to 8}{1 to 8}{0 to 7}{0 to 3}
\end{question}

\begin{question}{M}
In standard Ethernet a station gives up transmitting a frame after how many unsuccessful attempts?
\optsshort{16}{32}{10}{15}
\end{question}

\begin{question}{H}
Match List I with List II.\\
List I (IEEE standard): A. 802.3  B. 802.5  C. 802.11  D. 802.15\\
List II (Technology): I. Wireless LAN  II. Token ring  III. Bluetooth / wireless PAN  IV. CSMA/CD Ethernet
\optsshort{A-IV, B-I, C-II, D-III}{A-IV, B-II, C-III, D-I}{A-IV, B-II, C-I, D-III}{A-II, B-IV, C-I, D-III}
\end{question}

\begin{question}{M}
Match List I with List II.\\
List I (IEEE standard): A. 802.1Q  B. 802.2  C. 802.4  D. 802.16\\
List II (Technology): I. Token bus  II. WiMAX  III. VLAN tagging  IV. Logical link control
\optsshort{A-I, B-IV, C-III, D-II}{A-IV, B-III, C-I, D-II}{A-III, B-IV, C-II, D-I}{A-III, B-IV, C-I, D-II}
\end{question}

\begin{question}{H}
Which multiple access method is used by IEEE 802.11 wireless LANs?
\optsshort{CSMA/CA}{CSMA/CD}{Token passing}{Pure ALOHA}
\end{question}

\begin{question}{H}
Given below are two statements: one is labelled as Assertion A and the other is labelled as Reason R.\\
Assertion A: CSMA/CD is not used in wireless LANs.\\
Reason R: A wireless station cannot reliably detect a collision while transmitting, because its own signal is much stronger than any received signal and stations may be hidden from one another.
\optslong{Both A and R are true and R is the correct explanation of A}{Both A and R are true but R is NOT the correct explanation of A}{A is true but R is false}{A is false but R is true}
\end{question}

\begin{question}{H}
The RTS/CTS handshake in IEEE 802.11 is used to solve the
\optslong{silly window syndrome}{count-to-infinity problem}{exposed terminal problem only}{hidden terminal problem}
\end{question}

\begin{question}{H}
In the designation 10Base-T, the terms 10, Base and T stand for
\optslong{10 Mbps, broadband, twisted pair}{10 MHz, baseband, thick coax}{10 Mbps, baseband, twisted pair}{10 m, baseband, terminator}
\end{question}

\begin{question}{M}
The maximum segment length of 10Base5 (thick Ethernet) is
\optsshort{185 m}{500 m}{2000 m}{100 m}
\end{question}

\begin{question}{H}
An Ethernet MAC address is
\optslong{128 bits long and written in hexadecimal}{32 bits long and written in dotted decimal}{48 bits long and written in hexadecimal}{64 bits long and written in hexadecimal}
\end{question}

\begin{question}{M}
An Ethernet destination address in which the least significant bit of the first byte is 1 is a
\optsshort{multicast address}{loopback address}{private address}{unicast address}
\end{question}

\begin{question}{M}
The preamble and start frame delimiter of an IEEE 802.3 frame together occupy
\optsshort{12 bytes}{8 bytes}{7 bytes}{4 bytes}
\end{question}

\begin{question}{H}
A network has four hosts connected to a hub, the hub connected to a switch port, and three more hosts connected directly to other ports of the same switch. The number of collision domains and broadcast domains are respectively
\optsshort{1 and 1}{7 and 1}{4 and 1}{4 and 4}
\end{question}

\begin{question}{H}
Which device breaks up both collision domains and broadcast domains?
\optsshort{Router}{Bridge}{Layer-2 switch}{Hub}
\end{question}

\begin{question}{H}
A transparent (learning) bridge builds its forwarding table by examining the
\optslong{destination MAC address of incoming frames}{source MAC address of incoming frames}{destination IP address of incoming packets}{source IP address of incoming packets}
\end{question}

\begin{question}{M}
The spanning tree protocol is used in a switched LAN to
\optslong{find the shortest path between two hosts}{assign IP addresses to switches}{create VLANs}{prevent loops when redundant bridges are present}
\end{question}

\begin{question}{H}
Which of the following is NOT a type of HDLC frame?
\optsshort{Acknowledgement frame}{Unnumbered frame}{Supervisory frame}{Information frame}
\end{question}

\begin{question}{M}
The flag field that delimits an HDLC frame is
\optsshort{01111110}{10101010}{10101011}{11111111}
\end{question}

\begin{question}{M}
Which of the following statements about the Point-to-Point Protocol (PPP) is correct?
\optslong{PPP provides flow control using a sliding window}{PPP is a bit-oriented protocol using bit stuffing}{PPP can be used only on multipoint links}{PPP uses LCP to establish the link and NCP to configure the network layer}
\end{question}

\begin{question}{M}
The technique of temporarily delaying an acknowledgement so that it can be attached to the next outgoing data frame is called
\optsshort{bit stuffing}{piggybacking}{multiplexing}{pipelining}
\end{question}

\begin{question}{M}
Match List I with List II.\\
List I (Access method category): A. Random access  B. Controlled access  C. Channelization\\
List II (Protocols): I. FDMA, TDMA, CDMA  II. ALOHA, CSMA/CD, CSMA/CA  III. Reservation, polling, token passing
\optsshort{A-I, B-III, C-II}{A-III, B-II, C-I}{A-II, B-I, C-III}{A-II, B-III, C-I}
\end{question}

\begin{question}{M}
In which CSMA persistence method does a station that finds the channel idle transmit immediately with probability 1?
\optsshort{0-persistent}{1-persistent}{p-persistent}{Non-persistent}
\end{question}

\begin{question}{M}
Given below are two statements.\\
Statement I: In a datagram network every packet carries the full destination address and packets of the same message may follow different paths.\\
Statement II: In a virtual-circuit network resources are necessarily reserved for the entire duration of the connection, as in circuit switching.
\optslong{Both Statement I and Statement II are true}{Both Statement I and Statement II are false}{Statement I is true but Statement II is false}{Statement I is false but Statement II is true}
\end{question}

\begin{question}{M}
An ATM cell consists of
\optslong{5 bytes of header and 53 bytes of payload}{8 bytes of header and 53 bytes of payload}{48 bytes of header and 5 bytes of payload}{5 bytes of header and 48 bytes of payload}
\end{question}

\begin{question}{M}
Which line coding is used by IEEE 802.5 token ring?
\optsshort{Differential Manchester}{NRZ-L}{Manchester}{AMI}
\end{question}

\begin{question}{M}
Which of the following topologies requires a terminator at each end of the cable?
\optsshort{Star}{Bus}{Mesh}{Ring}
\end{question}

\begin{question}{M}
Given below are two statements.\\
Statement I: In a star topology the failure of the central hub brings down the whole network.\\
Statement II: In a star topology the failure of one link affects all the other stations.
\optslong{Both Statement I and Statement II are true}{Both Statement I and Statement II are false}{Statement I is true but Statement II is false}{Statement I is false but Statement II is true}
\end{question}

\begin{question}{H}
Arrange the layers of the OSI model from top to bottom.\\
A. Session  B. Network  C. Presentation  D. Transport  E. Data link
\optsshort{C, A, D, B, E}{A, C, D, B, E}{C, A, B, D, E}{C, D, A, B, E}
\end{question}

\begin{question}{H}
Which of the following are functions of the data link layer?\\
A. Framing  B. Flow control  C. Routing  D. Error control  E. Media access control\\
Choose the correct answer from the options given below:
\optsshort{A, B, D and E only}{A, B and C only}{B, C and D only}{A, C, D and E only}
\end{question}

\begin{question}{M}
In Go-Back-N ARQ, if frames 0 to 6 have been sent and the timer for frame 3 expires, the sender retransmits
\optsshort{frames 3, 4, 5 and 6}{frames 4, 5 and 6}{frame 3 only}{frames 0 to 6}
\end{question}

\begin{question}{M}
In Selective Repeat ARQ a negative acknowledgement NAK n means
\optslong{all frames from n onward are to be retransmitted}{all frames up to n are received}{the receiver window is full}{only frame n is to be retransmitted}
\end{question}

\begin{question}{H}
The link utilisation of Go-Back-N with window size N, when N < 1 + 2a, is
\optsshort{N / (1 + a)}{N / (1 + 2a)}{1 / (1 + 2a)}{1}
\end{question}

\begin{question}{M}
A LAN technology known as Fast Ethernet operates at
\optsshort{10 Gbps}{1 Gbps}{100 Mbps}{10 Mbps}
\end{question}

\begin{question}{M}
Which of the following networks covers the smallest geographical area?
\optsshort{PAN}{LAN}{MAN}{WAN}
\end{question}

\begin{question}{M}
The data link layer is divided by the IEEE 802 project into which two sublayers?
\optsshort{HDLC and PPP}{LLC and PHY}{MAC and PHY}{LLC and MAC}
\end{question}

\begin{question}{M}
In the byte-stuffing (character-stuffing) method of framing, a special byte is added to the data section when
\optslong{five consecutive 1s occur}{a parity error occurs}{the data contains a byte with the same pattern as the flag or the escape character}{the frame is too short}
\end{question}

\begin{question}{M}
Which of the following is a connection-oriented, packet-switched WAN technology that uses a 10-bit DLCI to identify virtual circuits?
\optsshort{Frame Relay}{Ethernet}{FDDI}{Token ring}
\end{question}

\begin{question}{M}
In IEEE 802.11, which interframe space gives the highest priority to frames such as ACK and CTS?
\optsshort{EIFS}{PIFS}{SIFS}{DIFS}
\end{question}

\begin{question}{M}
FDDI uses which of the following access methods and media?
\optslong{CSMA/CA over radio}{Polling over twisted pair}{Token passing over a dual fibre ring}{CSMA/CD over coaxial cable}
\end{question}

\begin{question}{M}
Given below are two statements.\\
Statement I: A VLAN divides a single switched network into multiple broadcast domains.\\
Statement II: Communication between hosts in different VLANs requires a layer-3 device.
\optslong{Both Statement I and Statement II are true}{Both Statement I and Statement II are false}{Statement I is true but Statement II is false}{Statement I is false but Statement II is true}
\end{question}

\begin{question}{H}
Given below are two statements: one is labelled as Assertion A and the other is labelled as Reason R.\\
Assertion A: Standard Ethernet specifies a minimum frame length of 64 bytes.\\
Reason R: A station must still be transmitting when a collision from the farthest station is detected, so the frame transmission time must be at least twice the maximum propagation time.
\optslong{Both A and R are true and R is the correct explanation of A}{Both A and R are true but R is NOT the correct explanation of A}{A is true but R is false}{A is false but R is true}
\end{question}

\begin{question}{M}
A group of computers and devices that share a common communication line in which control alternates between a primary station and secondary stations uses
\optsshort{reservation}{CSMA}{token passing}{polling}
\end{question}

\begin{question}{L}
Which of the following is a set of rules that governs data communication?
\optsshort{Topology}{Interface}{Protocol}{Standard}
\end{question}

\begin{question}{M}
The three key elements of a protocol are
\optslong{speed, cost and security}{sender, receiver and medium}{syntax, semantics and timing}{framing, routing and addressing}
\end{question}

\begin{question}{M}
Which of the following Ethernet standards uses optical fibre?
\optsshort{1000Base-T}{100Base-TX}{1000Base-SX}{10Base-T}
\end{question}

\begin{answers}
\ans{1}{2}{Physical: bits, data link: frames, network: packets (datagrams), transport: segments.}
\ans{2}{2}{Repeaters/hubs are layer 1, bridges/switches layer 2, routers layer 3, gateways can work up to the application layer.}
\ans{3}{3}{n(n $-$ 1)/2 = 10 $\times$ 9 / 2 = 45 duplex links.}
\ans{4}{2}{Each device connects directly to the other n $-$ 1 devices.}
\ans{5}{4}{The session layer establishes, manages and synchronises dialogs between applications.}
\ans{6}{1}{The presentation layer deals with the syntax and semantics of the information exchanged.}
\ans{7}{1}{Application data is wrapped into a transport segment, a network packet, a data link frame and finally sent as bits.}
\ans{8}{2}{TCP/IP has no separate session and presentation layers; their functions belong to the application layer.}
\ans{9}{2}{Port 16 bits, IPv4 32 bits, MAC 48 bits, IPv6 128 bits.}
\ans{10}{1}{a = Tp/Tt = 2; efficiency = 1/(1 + 2a) = 1/5 = 20\%.}
\ans{11}{2}{One frame per Tt + 2Tp = 1 + 9 = 10 ms, so 1000 bits / 10 ms = 100 kbps.}
\ans{12}{3}{Go-Back-N: 2\textsuperscript{m} $-$ 1 = 7; Selective Repeat: 2\textsuperscript{m$-$1} = 4.}
\ans{13}{3}{Window $\leq$ 2\textsuperscript{m$-$1}, so 2\textsuperscript{m$-$1} $\geq$ 32 gives m = 6.}
\ans{14}{3}{The Go-Back-N receiver accepts only the next in-order frame, so its window is 1.}
\ans{15}{3}{In one round-trip time the sender can transmit 1 + 2a frames.}
\ans{16}{4}{Pure ALOHA: S = Ge\textsuperscript{$-$2G}, max 1/(2e) at G = 0.5; slotted: S = Ge\textsuperscript{$-$G}, max 1/e at G = 1.}
\ans{17}{4}{Frame time is 1 ms, so G = 1; S = G e\textsuperscript{$-$2G} = 0.135, i.e. 135 frames per second.}
\ans{18}{2}{A frame can collide with frames started one frame time before or after it; in slotted ALOHA it is Tfr.}
\ans{19}{1}{Tt $\geq$ 2Tp = 2 $\times$ 12.5 $\mu$s = 25 $\mu$s; 25 $\mu$s $\times$ 10 Mbps = 250 bits.}
\ans{20}{2}{46--1500 bytes of payload plus 18 bytes of header and trailer gives 64--1518 bytes.}
\ans{21}{3}{After the n-th collision K is chosen uniformly from 0 to 2\textsuperscript{n} $-$ 1.}
\ans{22}{1}{The backoff exponent is capped at 10 and the frame is aborted after 16 attempts.}
\ans{23}{3}{802.3 Ethernet, 802.5 token ring, 802.11 WLAN, 802.15 WPAN.}
\ans{24}{4}{802.1Q VLAN, 802.2 LLC, 802.4 token bus, 802.16 WiMAX.}
\ans{25}{1}{Collisions cannot be detected reliably on a wireless medium, so 802.11 uses collision avoidance.}
\ans{26}{1}{Both the signal-strength problem and the hidden-station problem make collision detection impractical.}
\ans{27}{4}{RTS/CTS reserves the channel so that stations hidden from the sender hear the CTS and defer.}
\ans{28}{3}{Data rate in Mbps, baseband signalling, and twisted-pair medium.}
\ans{29}{2}{10Base5: 500 m; 10Base2: 185 m; 10Base-T: 100 m.}
\ans{30}{3}{A MAC address has 6 bytes (48 bits), usually written as 12 hexadecimal digits.}
\ans{31}{1}{LSB of the first byte 0 = unicast, 1 = multicast; all 1s is the broadcast address.}
\ans{32}{2}{7 bytes of preamble (alternating 1s and 0s) plus 1 byte of SFD (10101011).}
\ans{33}{3}{Each switch port is a collision domain (1 for the hub segment + 3 for the hosts); a switch does not divide the broadcast domain.}
\ans{34}{1}{Routers do not forward layer-2 broadcasts; switches and bridges split only collision domains.}
\ans{35}{2}{The bridge learns which port a station is on from the source address; it forwards using the destination address.}
\ans{36}{4}{STP blocks redundant ports to build a loop-free logical topology.}
\ans{37}{1}{HDLC defines I-frames, S-frames and U-frames; acknowledgements are carried in I- and S-frames.}
\ans{38}{1}{The flag 01111110 marks both the start and the end of a frame; bit stuffing keeps it unique.}
\ans{39}{4}{PPP is byte-oriented, uses byte stuffing, and has no flow control; LCP and NCP are its control protocols.}
\ans{40}{2}{Piggybacking saves bandwidth by carrying the ACK in the header of a data frame.}
\ans{41}{4}{Standard classification of multiple access protocols.}
\ans{42}{2}{1-persistent CSMA sends as soon as the line is idle, giving the highest chance of collision.}
\ans{43}{3}{A virtual circuit fixes the route and uses a VC identifier, but resource reservation is optional and may be on demand.}
\ans{44}{4}{An ATM cell is a fixed 53 bytes: 5-byte header plus 48-byte payload.}
\ans{45}{1}{Token ring uses differential Manchester; 10 Mbps Ethernet uses Manchester.}
\ans{46}{2}{Terminators absorb the signal at the ends of the bus backbone to prevent reflection.}
\ans{47}{3}{The hub is a single point of failure, but a broken link isolates only the station on that link.}
\ans{48}{1}{Application, Presentation, Session, Transport, Network, Data link, Physical.}
\ans{49}{1}{Routing is a network layer function; the others belong to the data link layer.}
\ans{50}{1}{Go-Back-N resends the lost frame and every frame sent after it.}
\ans{51}{4}{Selective Repeat retransmits only the damaged or lost frame.}
\ans{52}{2}{N frames are delivered in each cycle of duration (1 + 2a) frame times.}
\ans{53}{3}{Fast Ethernet (802.3u) is 100 Mbps; Gigabit Ethernet (802.3z/ab) is 1000 Mbps.}
\ans{54}{1}{A personal area network spans a few metres around one person.}
\ans{55}{4}{Logical link control (802.2) sits above the technology-specific media access control sublayer.}
\ans{56}{3}{An ESC byte is inserted before any flag or ESC pattern that appears in the data.}
\ans{57}{1}{Frame Relay identifies virtual circuits by the data link connection identifier.}
\ans{58}{3}{SIFS is the shortest interframe space, so frames waiting only SIFS get the medium first: SIFS < PIFS < DIFS.}
\ans{59}{3}{FDDI is a 100 Mbps token-passing network built as dual counter-rotating fibre rings.}
\ans{60}{1}{Each VLAN is its own broadcast domain, and inter-VLAN traffic must be routed.}
\ans{61}{1}{The minimum frame size comes directly from the CSMA/CD requirement Tt $\geq$ 2Tp.}
\ans{62}{4}{In polling the primary station selects and polls the secondaries.}
\ans{63}{3}{A protocol defines the syntax, semantics and timing of communication.}
\ans{64}{3}{Syntax (format), semantics (meaning) and timing (when and how fast).}
\ans{65}{3}{1000Base-SX is short-wavelength Gigabit Ethernet over multimode fibre; the others use twisted pair.}
\end{answers}

\section{Network Layer (IPv4/IPv6 addressing, subnetting, routing, congestion control)}

\begin{question}{H}
The IP address 191.200.10.5 belongs to
\optsshort{Class B}{Class C}{Class D}{Class A}
\end{question}

\begin{question}{H}
Match List I with List II.\\
List I (Class): A. Class A  B. Class B  C. Class C  D. Class D\\
List II (Leading bits of the address): I. 110  II. 1110  III. 0  IV. 10
\optsshort{A-III, B-IV, C-I, D-II}{A-III, B-I, C-IV, D-II}{A-IV, B-III, C-I, D-II}{A-III, B-IV, C-II, D-I}
\end{question}

\begin{question}{H}
The number of usable host addresses in a subnet with prefix /26 is
\optsshort{30}{62}{64}{126}
\end{question}

\begin{question}{H}
A host has the IP address 192.168.10.130/26. The network address and the broadcast address of its subnet are
\optslong{192.168.10.130 and 192.168.10.191}{192.168.10.128 and 192.168.10.255}{192.168.10.0 and 192.168.10.63}{192.168.10.128 and 192.168.10.191}
\end{question}

\begin{question}{H}
What is the subnet (network) address of the host 172.16.45.14/20?
\optsshort{172.16.0.0}{172.16.32.0}{172.16.45.0}{172.16.40.0}
\end{question}

\begin{question}{H}
The directed broadcast address of the subnet containing 172.16.45.14/20 is
\optsshort{172.16.47.255}{172.16.45.255}{172.16.63.255}{172.16.255.255}
\end{question}

\begin{question}{H}
A Class B network is subnetted with the mask 255.255.240.0. The number of subnets and the number of usable hosts per subnet are
\optsshort{16 and 4096}{16 and 4094}{240 and 4094}{4096 and 14}
\end{question}

\begin{question}{H}
An organisation needs a block that supports at least 500 hosts. The longest prefix (smallest block) that suffices is
\optsshort{/23}{/22}{/24}{/25}
\end{question}

\begin{question}{H}
A Class C network must be divided into 6 subnets. The subnet mask to be used and the number of usable hosts per subnet are
\optsshort{255.255.255.224 and 30}{255.255.255.224 and 32}{255.255.255.192 and 62}{255.255.255.240 and 14}
\end{question}

\begin{question}{H}
For the address 200.10.11.144/27, the last usable host address of the subnet is
\optsshort{200.10.11.175}{200.10.11.159}{200.10.11.190}{200.10.11.158}
\end{question}

\begin{question}{H}
The four networks 192.168.0.0/24, 192.168.1.0/24, 192.168.2.0/24 and 192.168.3.0/24 can be aggregated into the single supernet
\optsshort{192.168.0.0/23}{192.168.0.0/16}{192.168.0.0/22}{192.168.0.0/21}
\end{question}

\begin{question}{H}
How many addresses are contained in the CIDR block 20.10.30.32/28 and what is its last address?
\optsshort{16 and 20.10.30.47}{16 and 20.10.30.48}{32 and 20.10.30.63}{14 and 20.10.30.46}
\end{question}

\begin{question}{H}
Which of the following addresses belongs to the same subnet as 10.1.1.70/26?
\optsshort{10.1.1.60}{10.1.1.63}{10.1.1.130}{10.1.1.100}
\end{question}

\begin{question}{H}
Which of the following is NOT a private IPv4 address?
\optsshort{172.20.1.1}{192.168.100.1}{10.200.3.4}{172.32.1.1}
\end{question}

\begin{question}{M}
The address 127.0.0.1 is used for
\optsshort{multicasting}{limited broadcast}{loopback testing}{automatic private addressing}
\end{question}

\begin{question}{M}
A host that fails to obtain an address from a DHCP server assigns itself an address in the range
\optsshort{169.254.0.0/16}{224.0.0.0/4}{127.0.0.0/8}{10.0.0.0/8}
\end{question}

\begin{question}{M}
The number of networks in Class B and the number of addresses in each Class A network are respectively
\optsshort{2\textsuperscript{14} and 2\textsuperscript{16}}{2\textsuperscript{16} and 2\textsuperscript{24}}{2\textsuperscript{14} and 2\textsuperscript{24}}{2\textsuperscript{21} and 2\textsuperscript{8}}
\end{question}

\begin{question}{H}
A router has the following routing table entries: 10.0.0.0/8 $\rightarrow$ interface A, 10.1.0.0/16 $\rightarrow$ interface B, 10.1.2.0/24 $\rightarrow$ interface C, default $\rightarrow$ interface D. A packet destined to 10.1.3.9 is forwarded on
\optsshort{interface B}{interface A}{interface D}{interface C}
\end{question}

\begin{question}{H}
The minimum and maximum sizes of an IPv4 header are
\optsshort{20 bytes and 40 bytes}{8 bytes and 60 bytes}{20 bytes and 60 bytes}{40 bytes and 60 bytes}
\end{question}

\begin{question}{H}
An IPv4 datagram arrives with HLEN value 1000 (binary). The number of bytes of options carried is
\optsshort{32}{20}{8}{12}
\end{question}

\begin{question}{H}
An IPv4 datagram of total length 4000 bytes (20-byte header) has to pass through a link with an MTU of 1500 bytes. The number of fragments and the offset value of the last fragment are
\optsshort{3 and 2960}{2 and 185}{4 and 370}{3 and 370}
\end{question}

\begin{question}{H}
A fragment arrives with a fragmentation offset value of 100, HLEN of 5 and total length of 100. The numbers of its first and last data bytes are
\optsshort{100 and 179}{100 and 199}{800 and 879}{800 and 899}
\end{question}

\begin{question}{H}
In an IPv4 datagram the M (more fragments) bit is 0 and the fragmentation offset is non-zero. The datagram is
\optsshort{the first fragment}{the last fragment}{not fragmented}{a middle fragment}
\end{question}

\begin{question}{H}
Match List I with List II.\\
List I (Protocol): A. ICMP  B. TCP  C. UDP  D. OSPF\\
List II (Value of the IPv4 protocol field): I. 17  II. 89  III. 1  IV. 6
\optsshort{A-III, B-I, C-IV, D-II}{A-IV, B-III, C-I, D-II}{A-III, B-IV, C-II, D-I}{A-III, B-IV, C-I, D-II}
\end{question}

\begin{question}{M}
The purpose of the time-to-live (TTL) field in the IPv4 header is to
\optslong{measure the round-trip time}{indicate the age of the routing table}{order the fragments}{prevent a datagram from looping in the network forever}
\end{question}

\begin{question}{M}
The header checksum in an IPv4 datagram covers
\optslong{both header and data}{the header only}{the data only}{the header and a pseudo-header}
\end{question}

\begin{question}{H}
Which of the following fields of the IPv4 header is NOT present in the IPv6 base header?
\optsshort{Version}{Hop limit (TTL)}{Source address}{Header checksum}
\end{question}

\begin{question}{H}
The size of an IPv6 address and of the IPv6 base header are respectively
\optsshort{64 bits and 40 bytes}{128 bits and 20 bytes}{128 bytes and 40 bits}{128 bits and 40 bytes}
\end{question}

\begin{question}{H}
The IPv6 address 2001:0DB8:0000:0000:0000:FF00:0042:8329 can be abbreviated as
\optsshort{2001:DB8::FF::42:8329}{2001:DB8::FF00:42:8329}{2001:DB8:0:FF00:42:8329}{2001:DB8::FF:42:8329}
\end{question}

\begin{question}{H}
Which of the following types of address is NOT defined in IPv6?
\optsshort{Multicast}{Broadcast}{Unicast}{Anycast}
\end{question}

\begin{question}{M}
Given below are two statements.\\
Statement I: In IPv6, fragmentation can be performed only by the source host and not by intermediate routers.\\
Statement II: The IPv6 base header contains a 20-bit flow label field.
\optslong{Both Statement I and Statement II are true}{Both Statement I and Statement II are false}{Statement I is true but Statement II is false}{Statement I is false but Statement II is true}
\end{question}

\begin{question}{M}
Match List I with List II.\\
List I (IPv6 address/prefix): A. ::1  B. FE80::/10  C. FF00::/8  D. ::\\
List II (Meaning): I. Multicast  II. Unspecified address  III. Loopback  IV. Link-local unicast
\optsshort{A-IV, B-III, C-I, D-II}{A-III, B-IV, C-I, D-II}{A-III, B-I, C-IV, D-II}{A-II, B-IV, C-I, D-III}
\end{question}

\begin{question}{H}
Which protocol finds the physical (MAC) address corresponding to a known IP address?
\optsshort{DHCP}{ICMP}{RARP}{ARP}
\end{question}

\begin{question}{H}
Given below are two statements.\\
Statement I: An ARP request is broadcast on the local network.\\
Statement II: An ARP reply is also broadcast on the local network.
\optslong{Both Statement I and Statement II are true}{Both Statement I and Statement II are false}{Statement I is true but Statement II is false}{Statement I is false but Statement II is true}
\end{question}

\begin{question}{H}
Arrange the DHCP messages in the order in which they are exchanged when a host obtains an address.\\
A. DHCPREQUEST  B. DHCPOFFER  C. DHCPACK  D. DHCPDISCOVER
\optsshort{D, A, B, C}{B, D, A, C}{D, B, C, A}{D, B, A, C}
\end{question}

\begin{question}{H}
Which protocol is used by the ping and traceroute utilities?
\optsshort{ARP}{IGMP}{SNMP}{ICMP}
\end{question}

\begin{question}{M}
The traceroute program discovers the routers along a path by sending packets with
\optslong{the record-route option only}{increasing TTL values and collecting ICMP time-exceeded messages}{decreasing TTL values and collecting ICMP echo replies}{TCP SYN segments to every router}
\end{question}

\begin{question}{M}
Network address translation (NAT) is mainly used to
\optslong{allow many hosts with private addresses to share one or a few public addresses}{map IP addresses to MAC addresses}{translate domain names to IP addresses}{encrypt IP datagrams}
\end{question}

\begin{question}{H}
Match List I with List II.\\
List I (Routing protocol): A. RIP  B. OSPF  C. BGP\\
List II (Algorithm): I. Path vector  II. Distance vector  III. Link state
\optsshort{A-II, B-I, C-III}{A-III, B-II, C-I}{A-I, B-III, C-II}{A-II, B-III, C-I}
\end{question}

\begin{question}{H}
Match List I with List II.\\
List I (Routing protocol): A. RIP  B. OSPF  C. BGP\\
List II (Carried in): I. TCP port 179  II. IP directly, protocol number 89  III. UDP port 520
\optsshort{A-III, B-I, C-II}{A-III, B-II, C-I}{A-I, B-II, C-III}{A-II, B-III, C-I}
\end{question}

\begin{question}{H}
In RIP, a hop count of 16 denotes
\optslong{a directly connected network}{an unreachable destination (infinity)}{the default route}{the maximum allowed path}
\end{question}

\begin{question}{H}
The count-to-infinity problem is associated with
\optsshort{distance vector routing}{link state routing}{flooding}{path vector routing}
\end{question}

\begin{question}{H}
Distance vector routing and link state routing are based respectively on the
\optslong{Bellman-Ford and Dijkstra algorithms}{Dijkstra and Bellman-Ford algorithms}{Bellman-Ford and Floyd-Warshall algorithms}{Prim and Kruskal algorithms}
\end{question}

\begin{question}{M}
Given below are two statements.\\
Statement I: In distance vector routing each router shares its knowledge of the entire network only with its neighbours.\\
Statement II: In link state routing each router shares knowledge of its neighbourhood with every router in the area by flooding.
\optslong{Both Statement I and Statement II are true}{Both Statement I and Statement II are false}{Statement I is true but Statement II is false}{Statement I is false but Statement II is true}
\end{question}

\begin{question}{M}
Which of the following is an inter-domain (exterior gateway) routing protocol?
\optsshort{IS-IS}{BGP}{OSPF}{RIP}
\end{question}

\begin{question}{M}
RIP routers send periodic updates to their neighbours every
\optsshort{90 seconds}{10 seconds}{180 seconds}{30 seconds}
\end{question}

\begin{question}{M}
Which of the following routing techniques does not require any network information and always finds the shortest path, at the cost of generating many duplicate packets?
\optsshort{Flooding}{Link state routing}{Hierarchical routing}{Distance vector routing}
\end{question}

\begin{question}{H}
A token bucket has a capacity of 8 megabits and tokens arrive at 6 Mbps. If the maximum output rate is 10 Mbps and the bucket is initially full, for how long can the host transmit at the maximum rate?
\optsshort{2 s}{1.33 s}{0.8 s}{4 s}
\end{question}

\begin{question}{H}
Given below are two statements.\\
Statement I: The leaky bucket algorithm produces output at a constant rate regardless of the burstiness of the input.\\
Statement II: The token bucket algorithm allows an idle host to save up permission to send bursts later.
\optslong{Both Statement I and Statement II are true}{Both Statement I and Statement II are false}{Statement I is true but Statement II is false}{Statement I is false but Statement II is true}
\end{question}

\begin{question}{M}
Which of the following is an open-loop congestion control technique?
\optslong{Implicit signalling}{Backpressure}{Choke packet}{Admission/retransmission policy}
\end{question}

\begin{question}{M}
A choke packet is
\optslong{a packet with TTL equal to zero}{a packet sent by a congested router to the source asking it to reduce its rate}{a packet dropped by the router}{a corrupted packet}
\end{question}

\begin{question}{M}
IGMP is used for
\optslong{managing multicast group membership}{configuring hosts automatically}{resolving addresses}{reporting errors to the source}
\end{question}

\begin{question}{M}
Match List I with List II.\\
List I (ICMP message): A. Echo request  B. Destination unreachable  C. Time exceeded  D. Echo reply\\
List II (Type value): I. 0  II. 3  III. 11  IV. 8
\optsshort{A-IV, B-II, C-I, D-III}{A-IV, B-II, C-III, D-I}{A-I, B-II, C-III, D-IV}{A-IV, B-III, C-II, D-I}
\end{question}

\begin{question}{M}
The address 255.255.255.255 is
\optslong{a directed broadcast address}{a multicast address}{the limited broadcast address}{the loopback address}
\end{question}

\begin{question}{M}
Which strategies allow the coexistence of IPv4 and IPv6 during the transition?\\
A. Dual stack  B. Tunnelling  C. Header translation  D. Supernetting\\
Choose the correct answer from the options given below:
\optsshort{A, B and C only}{A, B, C and D}{A and D only}{B, C and D only}
\end{question}

\begin{question}{M}
In classful addressing, the default mask of a Class B address is
\optsshort{255.0.0.0}{255.255.255.0}{255.255.0.0}{255.255.255.255}
\end{question}

\begin{question}{M}
The identification, flags and fragmentation offset fields of the IPv4 header have sizes (in bits) of
\optsshort{32, 3 and 13}{8, 3 and 13}{16, 4 and 12}{16, 3 and 13}
\end{question}

\begin{question}{H}
An ISP is granted the block 16.12.64.0/20. The total number of addresses in the block is
\optsshort{4096}{2048}{8192}{1024}
\end{question}

\begin{question}{H}
The subnet mask 255.255.255.248 in a Class C network provides
\optslong{16 subnets with 14 hosts each}{32 subnets with 6 hosts each}{32 subnets with 8 hosts each}{8 subnets with 30 hosts each}
\end{question}

\begin{question}{M}
Which of the following services is provided by the IPv4 protocol?
\optslong{Reliable, connection-oriented delivery}{Flow-controlled delivery}{Unreliable, connectionless best-effort delivery}{Reliable, connectionless delivery}
\end{question}

\begin{question}{M}
In OSPF, an autonomous system is divided into areas, all of which must be connected to
\optsshort{a transit AS}{area 1}{the backbone area (area 0)}{the stub area}
\end{question}

\begin{question}{M}
Which of the following protocols supports quality of service by reserving resources along a path for a flow?
\optsshort{RIP}{ICMP}{RSVP}{ARP}
\end{question}

\begin{question}{M}
Split horizon is a technique in which a router
\optslong{sends the full table to all neighbours}{uses two metrics for each route}{does not advertise a route back on the interface from which it was learned}{floods link state packets}
\end{question}

\begin{question}{M}
Which of the following is a multicast (Class D) address?
\optsshort{240.0.0.1}{192.0.2.1}{172.16.0.1}{224.0.0.5}
\end{question}

\begin{question}{H}
Given below are two statements: one is labelled as Assertion A and the other is labelled as Reason R.\\
Assertion A: The IPv4 header checksum must be recalculated at every router.\\
Reason R: The TTL field of the header changes at every hop.
\optslong{Both A and R are true and R is the correct explanation of A}{Both A and R are true but R is NOT the correct explanation of A}{A is true but R is false}{A is false but R is true}
\end{question}

\begin{question}{M}
The maximum total length of an IPv4 datagram is
\optsshort{1500 bytes}{4,294,967,295 bytes}{65,535 bytes}{65,536 bytes}
\end{question}

\begin{question}{M}
Which of the following addresses can be assigned to a host on the subnet 192.168.5.32/27?
\optsshort{192.168.5.63}{192.168.5.50}{192.168.5.64}{192.168.5.32}
\end{question}

\begin{question}{M}
A router performs the bitwise AND of a destination address with the subnet mask in order to obtain the
\optsshort{subnetwork address}{broadcast address}{next-hop MAC address}{host address}
\end{question}

\begin{answers}
\ans{1}{1}{Class B covers first octets 128--191.}
\ans{2}{1}{A: 0, B: 10, C: 110, D: 1110, E: 1111.}
\ans{3}{2}{6 host bits give 2\textsuperscript{6} $-$ 2 = 62 usable addresses.}
\ans{4}{4}{A /26 has blocks of 64; 130 lies in the block 128--191.}
\ans{5}{2}{Mask 255.255.240.0; 45 AND 240 = 32, so the subnet is 172.16.32.0.}
\ans{6}{1}{The block is 172.16.32.0 -- 172.16.47.255.}
\ans{7}{2}{4 bits are borrowed (16 subnets), leaving 12 host bits: 2\textsuperscript{12} $-$ 2 = 4094.}
\ans{8}{1}{/23 gives 2\textsuperscript{9} $-$ 2 = 510 usable addresses; /24 gives only 254.}
\ans{9}{1}{3 bits are needed for 6 subnets (2\textsuperscript{3} = 8), leaving 5 host bits: 30 hosts.}
\ans{10}{4}{Blocks of 32: 128--159; the broadcast is .159, so the last host is .158.}
\ans{11}{3}{Four contiguous /24 blocks share the first 22 bits.}
\ans{12}{1}{32 $-$ 28 = 4 host bits give 16 addresses: .32 to .47.}
\ans{13}{4}{The /26 block containing .70 is .64 to .127.}
\ans{14}{4}{Private ranges: 10.0.0.0/8, 172.16.0.0/12 (172.16--172.31) and 192.168.0.0/16.}
\ans{15}{3}{The 127.0.0.0/8 block is reserved for loopback.}
\ans{16}{1}{169.254.0.0/16 is the link-local (APIPA) block.}
\ans{17}{3}{Class B has a 16-bit netid with 2 fixed bits (2\textsuperscript{14} networks); Class A has 24 host bits.}
\ans{18}{1}{Longest prefix match: /8 and /16 match, /24 does not; the /16 entry is chosen.}
\ans{19}{3}{HLEN is 4 bits counted in 4-byte words: 5 $\times$ 4 = 20 to 15 $\times$ 4 = 60 bytes.}
\ans{20}{4}{HLEN = 8, so the header is 32 bytes; 32 $-$ 20 = 12 bytes of options.}
\ans{21}{4}{3980 data bytes split as 1480 + 1480 + 1020; offsets are 0, 185 and 370 (in units of 8 bytes).}
\ans{22}{3}{First byte = 100 $\times$ 8 = 800; data length = 100 $-$ 20 = 80 bytes, so the last byte is 879.}
\ans{23}{2}{M = 0 with offset $\neq$ 0 means last fragment; M = 0 with offset 0 means unfragmented.}
\ans{24}{4}{ICMP 1, TCP 6, UDP 17, OSPF 89.}
\ans{25}{4}{Each router decrements TTL; the datagram is discarded when it reaches 0.}
\ans{26}{2}{IPv4 checks only its header, and the checksum is recomputed at every router because TTL changes.}
\ans{27}{4}{IPv6 drops the header checksum, HLEN and the fragmentation fields from the base header.}
\ans{28}{4}{16-byte addresses and a fixed 40-byte base header.}
\ans{29}{2}{Leading zeros of each group are dropped and one run of all-zero groups is replaced by ::.}
\ans{30}{2}{IPv6 replaces broadcast with multicast and adds anycast.}
\ans{31}{1}{IPv6 routers never fragment (a fragment extension header is used by the source), and the flow label is 20 bits.}
\ans{32}{2}{::1 loopback, FE80::/10 link-local, FF00::/8 multicast, :: unspecified.}
\ans{33}{4}{ARP maps a logical address to a physical address; RARP does the reverse.}
\ans{34}{3}{The reply is unicast directly to the requester, whose MAC address is in the request.}
\ans{35}{4}{Discover, Offer, Request, Acknowledge (DORA).}
\ans{36}{4}{Ping uses ICMP echo request/reply; traceroute relies on ICMP time-exceeded messages.}
\ans{37}{2}{Each successive probe has a TTL one larger, so each router in turn discards it and reports back.}
\ans{38}{1}{NAT rewrites addresses (and ports) at the border router, conserving public IPv4 addresses.}
\ans{39}{4}{RIP is distance vector, OSPF is link state, BGP is path vector.}
\ans{40}{2}{RIP uses UDP 520, OSPF runs directly over IP (89), BGP uses TCP 179.}
\ans{41}{2}{RIP limits paths to 15 hops; 16 means infinity.}
\ans{42}{1}{Bad news propagates slowly in distance vector routing; split horizon and poison reverse reduce the problem.}
\ans{43}{1}{Distance vector uses the Bellman-Ford equation; link state runs Dijkstra on the full topology.}
\ans{44}{1}{This is the classic contrast between the two approaches.}
\ans{45}{2}{BGP routes between autonomous systems; the others are intra-domain protocols.}
\ans{46}{4}{The RIP periodic timer is about 30 s; the expiration timer is 180 s.}
\ans{47}{1}{Flooding sends every packet on every outgoing line, so one copy necessarily takes the shortest path.}
\ans{48}{1}{S = C / (M $-$ $\rho$) = 8 / (10 $-$ 6) = 2 s.}
\ans{49}{1}{Leaky bucket smooths traffic rigidly; token bucket permits bounded bursts.}
\ans{50}{4}{Open-loop methods prevent congestion through policies; choke packets, backpressure and signalling are closed-loop.}
\ans{51}{2}{It is an explicit closed-loop congestion signal directed at the source.}
\ans{52}{1}{Hosts use IGMP to inform local routers of the multicast groups they wish to join.}
\ans{53}{2}{Echo reply 0, destination unreachable 3, echo request 8, time exceeded 11.}
\ans{54}{3}{It reaches all hosts on the local network only; routers never forward it.}
\ans{55}{1}{Supernetting is an IPv4 address aggregation technique, not a transition strategy.}
\ans{56}{3}{Class B has a 16-bit network part.}
\ans{57}{4}{Together they make up one 32-bit word of the header.}
\ans{58}{1}{2\textsuperscript{32 $-$ 20} = 4096 addresses.}
\ans{59}{2}{248 = 11111000: 5 subnet bits (32 subnets) and 3 host bits (6 hosts).}
\ans{60}{3}{IP is a best-effort datagram protocol; reliability is left to upper layers.}
\ans{61}{3}{All OSPF areas attach to the backbone, area 0, through area border routers.}
\ans{62}{3}{RSVP is the signalling protocol of Integrated Services.}
\ans{63}{3}{It helps avoid two-node routing loops in distance vector protocols.}
\ans{64}{4}{Class D spans 224.0.0.0 to 239.255.255.255; 240 and above is Class E.}
\ans{65}{1}{Any change in the header (TTL, and possibly fragmentation fields) invalidates the previous checksum.}
\ans{66}{3}{The total length field is 16 bits, so the maximum is 2\textsuperscript{16} $-$ 1 bytes.}
\ans{67}{2}{.32 is the network address and .63 the broadcast address; hosts are .33 to .62.}
\ans{68}{1}{Masking extracts the network/subnet prefix used for the routing table lookup.}
\end{answers}

\section{Transport Layer, Application Layer Protocols and the World Wide Web}

\begin{question}{H}
The sizes of the UDP header and the minimum TCP header are respectively
\optsshort{20 bytes and 8 bytes}{8 bytes and 40 bytes}{20 bytes and 20 bytes}{8 bytes and 20 bytes}
\end{question}

\begin{question}{H}
Match List I with List II.\\
List I (Protocol): A. FTP (control)  B. SMTP  C. DNS  D. HTTPS\\
List II (Well-known port): I. 53  II. 443  III. 21  IV. 25
\optsshort{A-III, B-I, C-IV, D-II}{A-III, B-IV, C-I, D-II}{A-III, B-IV, C-II, D-I}{A-IV, B-III, C-I, D-II}
\end{question}

\begin{question}{H}
Match List I with List II.\\
List I (Protocol): A. Telnet  B. SSH  C. POP3  D. IMAP\\
List II (Well-known port): I. 143  II. 110  III. 22  IV. 23
\optsshort{A-IV, B-III, C-II, D-I}{A-II, B-III, C-IV, D-I}{A-IV, B-III, C-I, D-II}{A-III, B-IV, C-II, D-I}
\end{question}

\begin{question}{H}
Which of the following application layer protocols normally use UDP as the transport protocol?\\
A. DNS  B. SMTP  C. TFTP  D. SNMP  E. HTTP/1.1\\
Choose the correct answer from the options given below:
\optsshort{B and E only}{A, C, D and E only}{A, C and D only}{A, B and C only}
\end{question}

\begin{question}{H}
The range of well-known port numbers is
\optsshort{0 to 255}{1024 to 49151}{49152 to 65535}{0 to 1023}
\end{question}

\begin{question}{H}
Arrange the segments exchanged during TCP connection establishment in the correct order.\\
A. ACK  B. SYN  C. SYN + ACK
\optsshort{B, C, A}{A, B, C}{B, A, C}{C, B, A}
\end{question}

\begin{question}{H}
A TCP receiver gets a segment with sequence number 5000 carrying 1000 bytes of data, and all earlier bytes have been received. The acknowledgement number it sends is
\optsshort{5000}{5999}{6001}{6000}
\end{question}

\begin{question}{H}
A TCP connection starts in slow start with a congestion window of 1 MSS and a slow-start threshold of 8 MSS. If no loss occurs, the congestion window during the 6th transmission round is
\optsshort{16 MSS}{12 MSS}{32 MSS}{10 MSS}
\end{question}

\begin{question}{H}
A TCP Tahoe sender has a congestion window of 24 MSS when a timeout occurs. The new values of the slow-start threshold and the congestion window are
\optsshort{12 MSS and 1 MSS}{12 MSS and 12 MSS}{12 MSS and 15 MSS}{24 MSS and 1 MSS}
\end{question}

\begin{question}{H}
In TCP Reno, on receiving three duplicate acknowledgements the sender
\optslong{closes the connection}{sets the congestion window to 1 MSS and enters slow start}{halves the congestion window and continues with congestion avoidance (fast recovery)}{doubles the congestion window}
\end{question}

\begin{question}{H}
TCP's congestion control in the congestion avoidance phase follows
\optslong{multiplicative increase, additive decrease}{multiplicative increase, multiplicative decrease}{additive increase, additive decrease}{additive increase, multiplicative decrease}
\end{question}

\begin{question}{H}
The actual size of the TCP sender's window is
\optslong{the maximum of the congestion window and the receiver's advertised window}{the receiver's advertised window}{the minimum of the congestion window and the receiver's advertised window}{the congestion window}
\end{question}

\begin{question}{M}
The estimated RTT of a TCP connection is 100 ms. A new RTT sample of 140 ms is measured. With $\alpha$ = 0.125, the new estimated RTT is
\optsshort{135 ms}{105 ms}{110 ms}{120 ms}
\end{question}

\begin{question}{M}
Match List I with List II.\\
List I (Algorithm): A. Nagle's algorithm  B. Clark's solution  C. Karn's algorithm\\
List II (Purpose): I. Receiver-side remedy for silly window syndrome  II. Excludes retransmitted segments from RTT estimation  III. Sender-side remedy for silly window syndrome
\optsshort{A-II, B-I, C-III}{A-III, B-II, C-I}{A-I, B-III, C-II}{A-III, B-I, C-II}
\end{question}

\begin{question}{M}
The sizes of the sequence number field and the window size field in the TCP header are
\optsshort{16 bits and 16 bits}{16 bits and 32 bits}{32 bits and 32 bits}{32 bits and 16 bits}
\end{question}

\begin{question}{M}
Which TCP flag is used to abort a connection abruptly?
\optsshort{PSH}{URG}{RST}{FIN}
\end{question}

\begin{question}{M}
A socket address is the combination of
\optslong{a port number and a protocol number}{a MAC address and an IP address}{a MAC address and a port number}{an IP address and a port number}
\end{question}

\begin{question}{M}
Given below are two statements.\\
Statement I: TCP is a byte-stream-oriented protocol that numbers every byte it transmits.\\
Statement II: UDP provides flow control and congestion control but no error control.
\optslong{Both Statement I and Statement II are true}{Both Statement I and Statement II are false}{Statement I is true but Statement II is false}{Statement I is false but Statement II is true}
\end{question}

\begin{question}{M}
The pseudo-header used in the computation of the TCP and UDP checksums contains
\optslong{the source and destination port numbers}{the source and destination IP addresses, the protocol number and the segment length}{the sequence and acknowledgement numbers}{the MAC addresses of sender and receiver}
\end{question}

\begin{question}{M}
A SYN flooding attack exploits
\optslong{the fragmentation mechanism of IP}{the three-way handshake of TCP}{the checksum of UDP}{the four-way connection termination of TCP}
\end{question}

\begin{question}{M}
Which of the following features is provided by SCTP but not by TCP?
\optslong{Reliable delivery}{Multihoming and multiple streams per association}{Connection-oriented service}{Congestion control}
\end{question}

\begin{question}{H}
Which of the following statements about HTTP is correct?
\optslong{HTTP is a stateful protocol that uses UDP}{HTTP is a push protocol}{HTTP uses two parallel TCP connections for control and data}{HTTP is a stateless protocol that uses TCP port 80}
\end{question}

\begin{question}{H}
A web page consists of a base HTML file and 10 embedded images on the same server. Using non-persistent HTTP without parallel connections, how many TCP connections are needed?
\optsshort{1}{2}{10}{11}
\end{question}

\begin{question}{H}
In non-persistent HTTP, the time needed to fetch one small object (ignoring transmission time), where RTT is the round-trip time, is approximately
\optsshort{2 RTT}{3 RTT}{4 RTT}{1 RTT}
\end{question}

\begin{question}{H}
Match List I with List II.\\
List I (HTTP status code): A. 200  B. 301  C. 404  D. 500\\
List II (Meaning): I. Not found  II. Internal server error  III. OK  IV. Moved permanently
\optsshort{A-IV, B-III, C-I, D-II}{A-III, B-IV, C-I, D-II}{A-III, B-I, C-IV, D-II}{A-III, B-IV, C-II, D-I}
\end{question}

\begin{question}{M}
Which HTTP method retrieves only the headers of a resource, without the body?
\optsshort{PUT}{HEAD}{GET}{POST}
\end{question}

\begin{question}{M}
Which mechanism lets a web server recognise a returning user even though HTTP is stateless?
\optsshort{Three-way handshake}{DNS caching}{MIME}{Cookies}
\end{question}

\begin{question}{H}
Which of the following statements about FTP is correct?
\optslong{FTP uses a single connection for both control and data}{FTP runs over UDP port 69}{FTP uses TCP port 21 for control and TCP port 20 for data}{FTP is a stateless protocol}
\end{question}

\begin{question}{H}
Given below are two statements.\\
Statement I: SMTP is a push protocol used to transfer mail from the sender's mail server to the receiver's mail server.\\
Statement II: POP3 and IMAP are pull protocols used by a user agent to retrieve mail from the mail server.
\optslong{Both Statement I and Statement II are true}{Both Statement I and Statement II are false}{Statement I is true but Statement II is false}{Statement I is false but Statement II is true}
\end{question}

\begin{question}{M}
Which protocol allows non-ASCII data such as images and audio to be sent through e-mail?
\optsshort{ICMP}{SNMP}{POP3}{MIME}
\end{question}

\begin{question}{H}
Match List I with List II.\\
List I (DNS resource record): A. A  B. AAAA  C. MX  D. CNAME\\
List II (Meaning): I. Mail exchanger for the domain  II. Alias (canonical name) mapping  III. IPv4 address of a host  IV. IPv6 address of a host
\optsshort{A-IV, B-III, C-I, D-II}{A-III, B-IV, C-I, D-II}{A-I, B-IV, C-III, D-II}{A-III, B-IV, C-II, D-I}
\end{question}

\begin{question}{H}
The Domain Name System normally uses
\optslong{UDP port 67}{TCP port 53 for all messages}{TCP port 80}{UDP port 53 for queries and TCP port 53 for zone transfers and large responses}
\end{question}

\begin{question}{M}
In a recursive DNS query
\optslong{the queried server takes responsibility for resolving the name completely and returns the final answer}{the client contacts each name server in turn}{no caching is done}{only root servers are contacted}
\end{question}

\begin{question}{M}
Which DNS record type is used for reverse lookup, i.e. mapping an IP address to a name?
\optsshort{NS}{SOA}{PTR}{TXT}
\end{question}

\begin{question}{M}
In SNMP, the collection of managed objects maintained by an agent is called the
\optsshort{ASN}{SMI}{PDU}{MIB}
\end{question}

\begin{question}{M}
SNMP agents listen on UDP port 161, and managers receive trap messages on UDP port
\optsshort{160}{514}{163}{162}
\end{question}

\begin{question}{M}
Which of the following provides secure remote login and has replaced Telnet?
\optsshort{FTP}{SNMP}{SSH}{TFTP}
\end{question}

\begin{question}{H}
Arrange the following system calls in the order in which a TCP server invokes them.\\
A. listen()  B. socket()  C. accept()  D. bind()
\optsshort{B, A, D, C}{B, D, C, A}{B, D, A, C}{D, B, A, C}
\end{question}

\begin{question}{M}
The World Wide Web was invented by
\optslong{Vint Cerf at DARPA}{Tim Berners-Lee at CERN}{Robert Metcalfe at Xerox PARC}{Ray Tomlinson at BBN}
\end{question}

\begin{question}{M}
In the URL http://www.example.org:8080/docs/index.html, the component 8080 is the
\optsshort{host}{protocol}{port number}{path}
\end{question}

\begin{question}{M}
Match List I with List II.\\
List I (Type of web document): A. Static document  B. Dynamic document  C. Active document\\
List II (Description): I. A program run at the client site, e.g. a Java applet or script  II. Fixed-content file stored on the server  III. Created by the server on each request, e.g. using CGI
\optsshort{A-II, B-I, C-III}{A-II, B-III, C-I}{A-III, B-II, C-I}{A-I, B-III, C-II}
\end{question}

\begin{question}{M}
A proxy server (web cache) improves performance by
\optslong{compressing the TCP headers}{keeping copies of recently requested objects and serving them locally}{assigning IP addresses dynamically}{encrypting the HTTP messages}
\end{question}

\begin{question}{M}
Which HTTP request header and status code are used by a cache to verify that its stored copy is still up to date?
\optslong{Range and 404 Not Found}{Host and 200 OK}{Cookie and 401 Unauthorized}{If-Modified-Since and 304 Not Modified}
\end{question}

\begin{question}{M}
Given below are two statements.\\
Statement I: HTTP/1.1 uses persistent connections by default.\\
Statement II: HTTPS is HTTP running over TLS/SSL and uses port 443 by default.
\optslong{Both Statement I and Statement II are true}{Both Statement I and Statement II are false}{Statement I is true but Statement II is false}{Statement I is false but Statement II is true}
\end{question}

\begin{question}{M}
Which protocol is used to synchronise the clocks of computers over a network?
\optsshort{RTP}{NTP}{SNMP}{NNTP}
\end{question}

\begin{question}{M}
Real-time Transport Protocol (RTP), used for audio and video streaming, normally runs over
\optsshort{UDP}{ICMP}{ARP}{TCP}
\end{question}

\begin{question}{M}
TFTP differs from FTP in that TFTP
\optslong{uses UDP port 69 and provides no authentication}{supports directory listing}{uses two connections}{uses TCP port 21}
\end{question}

\begin{question}{M}
Which of the following correctly orders the steps taken when a user types a URL into a browser for the first time?\\
A. TCP connection establishment  B. DNS lookup of the host name  C. HTTP request  D. HTTP response
\optsshort{C, B, A, D}{A, B, C, D}{B, A, C, D}{B, C, A, D}
\end{question}

\begin{question}{H}
Given below are two statements: one is labelled as Assertion A and the other is labelled as Reason R.\\
Assertion A: DNS uses UDP for ordinary name queries.\\
Reason R: A query and its response are short, and avoiding connection set-up reduces the delay.
\optslong{Both A and R are true and R is the correct explanation of A}{Both A and R are true but R is NOT the correct explanation of A}{A is true but R is false}{A is false but R is true}
\end{question}

\begin{question}{M}
In TCP, each of the SYN and FIN segments
\optslong{consumes one sequence number}{consumes as many sequence numbers as the window size}{consumes no sequence number}{cannot carry an acknowledgement}
\end{question}

\begin{question}{M}
The maximum size of a TCP receive window without the window scaling option is
\optsshort{1024 bytes}{4 GB}{1 MB}{65,535 bytes}
\end{question}

\begin{question}{M}
Which transport-layer function allows several application processes on one host to share the network using port numbers?
\optslong{Multiplexing and demultiplexing}{Flow control}{Congestion control}{Segmentation}
\end{question}

\begin{question}{M}
In IMAP, unlike POP3
\optslong{mail is sent from client to server}{mail is always deleted from the server after download}{UDP is used}{a user can organise mail in folders on the server and download only headers}
\end{question}

\begin{question}{M}
Which of the following is a markup language used to describe the structure and transport of data rather than its presentation?
\optsshort{CSS}{JavaScript}{XML}{HTML}
\end{question}

\begin{question}{M}
The HTTP status code 403 indicates
\optsshort{Service unavailable}{Forbidden}{Unauthorized}{Bad request}
\end{question}

\begin{question}{L}
HTTP/3 runs over which transport protocol?
\optsshort{QUIC over UDP}{SCTP}{DCCP}{TCP}
\end{question}

\begin{question}{M}
Which of the following is a peer-to-peer file sharing protocol in which a file is divided into chunks exchanged among a swarm of peers?
\optsshort{SMTP}{BitTorrent}{NFS}{FTP}
\end{question}

\begin{question}{M}
Which of the following is NOT a service provided by TCP?
\optsshort{Full-duplex communication}{Guaranteed minimum bandwidth}{Flow control}{Reliable in-order delivery}
\end{question}

\begin{answers}
\ans{1}{4}{UDP has a fixed 8-byte header; TCP's header is 20 to 60 bytes.}
\ans{2}{2}{FTP control 21, SMTP 25, DNS 53, HTTPS 443.}
\ans{3}{1}{Telnet 23, SSH 22, POP3 110, IMAP 143.}
\ans{4}{3}{DNS, TFTP and SNMP run over UDP; SMTP and HTTP/1.1 use TCP.}
\ans{5}{4}{Well-known 0--1023, registered 1024--49151, dynamic/ephemeral 49152--65535.}
\ans{6}{1}{Three-way handshake: SYN, SYN+ACK, ACK.}
\ans{7}{4}{TCP acknowledges the next byte expected: 5000 + 1000 = 6000.}
\ans{8}{4}{Rounds: 1, 2, 4, 8 (threshold reached), then additive increase: 9, 10.}
\ans{9}{1}{On timeout ssthresh becomes half the current window and cwnd restarts from 1 MSS.}
\ans{10}{3}{Reno performs fast retransmit followed by fast recovery instead of slow start.}
\ans{11}{4}{AIMD: the window grows by 1 MSS per RTT and is halved on loss.}
\ans{12}{3}{The sender must respect both flow control (rwnd) and congestion control (cwnd).}
\ans{13}{2}{EstimatedRTT = 0.875 $\times$ 100 + 0.125 $\times$ 140 = 105 ms.}
\ans{14}{4}{Nagle acts at the sender, Clark at the receiver; Karn's algorithm handles the retransmission ambiguity.}
\ans{15}{4}{Sequence and acknowledgement numbers are 32 bits; the advertised window is 16 bits.}
\ans{16}{3}{RST resets the connection; FIN performs a graceful close.}
\ans{17}{4}{IP address identifies the host and the port identifies the process.}
\ans{18}{3}{UDP offers no flow or congestion control; it has only an optional checksum.}
\ans{19}{2}{The pseudo-header guards against misdelivered datagrams.}
\ans{20}{2}{The attacker sends many SYNs with spoofed addresses, exhausting the server's half-open connection table.}
\ans{21}{2}{SCTP is message-oriented and supports multistreaming and multihoming.}
\ans{22}{4}{HTTP itself keeps no client state; cookies are used to add state.}
\ans{23}{4}{Non-persistent HTTP opens a separate connection for each of the 11 objects.}
\ans{24}{1}{One RTT to set up the TCP connection and one RTT for the request and response.}
\ans{25}{2}{2xx success, 3xx redirection, 4xx client error, 5xx server error.}
\ans{26}{2}{HEAD is identical to GET except that the server does not return a message body.}
\ans{27}{4}{A cookie stored by the browser is sent back with each later request to the same server.}
\ans{28}{3}{FTP sends control information out-of-band on a separate connection and maintains user state.}
\ans{29}{1}{SMTP pushes mail; message access agents pull it.}
\ans{30}{4}{MIME converts non-ASCII data to 7-bit ASCII (e.g. Base64) for SMTP.}
\ans{31}{2}{A: IPv4, AAAA: IPv6, MX: mail server, CNAME: alias.}
\ans{32}{4}{Ordinary queries use UDP; responses that do not fit and zone transfers use TCP.}
\ans{33}{1}{In an iterative query the server instead returns a referral to another server.}
\ans{34}{3}{PTR records in the in-addr.arpa domain support inverse resolution.}
\ans{35}{4}{The Management Information Base is the agent's virtual database; SMI defines the rules for naming objects.}
\ans{36}{4}{Port 161 for requests to agents, 162 for traps sent to the manager.}
\ans{37}{3}{SSH encrypts the session; Telnet sends everything, including passwords, in plain text.}
\ans{38}{3}{socket, bind, listen, accept.}
\ans{39}{2}{Berners-Lee proposed the Web at CERN in 1989.}
\ans{40}{3}{A URL has the form protocol://host:port/path.}
\ans{41}{2}{Static: fixed; dynamic: generated at the server; active: executed at the client.}
\ans{42}{2}{A cache hit avoids contacting the origin server, reducing response time and traffic.}
\ans{43}{4}{A conditional GET returns 304 with no body if the object is unchanged.}
\ans{44}{1}{Both statements are true.}
\ans{45}{2}{The Network Time Protocol runs on UDP port 123.}
\ans{46}{1}{Timeliness matters more than reliability for real-time media, so RTP uses UDP.}
\ans{47}{1}{TFTP is a minimal file transfer protocol over UDP.}
\ans{48}{3}{The name is resolved first, then the connection is set up, then request and response.}
\ans{49}{1}{The low overhead of UDP suits short request-response exchanges.}
\ans{50}{1}{SYN and FIN each consume one sequence number even though they carry no data.}
\ans{51}{4}{The window field is 16 bits wide.}
\ans{52}{1}{Ports let the transport layer multiplex at the sender and demultiplex at the receiver.}
\ans{53}{4}{IMAP keeps mail on the server and supports folders and partial fetch.}
\ans{54}{3}{XML defines user-defined tags for describing data; HTML describes how to display it.}
\ans{55}{2}{400 Bad Request, 401 Unauthorized, 403 Forbidden, 503 Service Unavailable.}
\ans{56}{1}{HTTP/3 replaces TCP with QUIC, which is built on UDP.}
\ans{57}{2}{BitTorrent peers trade chunks using a tit-for-tat strategy, coordinated by a tracker or DHT.}
\ans{58}{2}{TCP offers no timing or bandwidth guarantees.}
\end{answers}

\section{Network Security (malware, cryptography, steganography, digital signatures, VPNs, firewalls)}

\begin{question}{H}
In RSA, p = 3, q = 11 and the public exponent e = 7. The private exponent d is
\optsshort{7}{3}{17}{13}
\end{question}

\begin{question}{H}
Using RSA with p = 3, q = 11 and e = 7, the ciphertext for plaintext M = 2 is
\optsshort{29}{8}{17}{31}
\end{question}

\begin{question}{H}
In an RSA system p = 7, q = 11 and e = 13. The value of d is
\optsshort{17}{37}{43}{23}
\end{question}

\begin{question}{H}
In RSA with p = 5 and q = 11, the public key is (e, n) = (3, 55). The private key d is
\optsshort{27}{7}{37}{13}
\end{question}

\begin{question}{H}
In the Diffie-Hellman key exchange the prime is q = 23 and the primitive root is $\alpha$ = 5. If user A selects the private key 6, A's public key is
\optsshort{2}{15}{19}{8}
\end{question}

\begin{question}{H}
In a Diffie-Hellman exchange with q = 23 and $\alpha$ = 5, A's private key is 6 and B's private key is 15. The shared secret key is
\optsshort{8}{19}{2}{13}
\end{question}

\begin{question}{H}
How many keys are required for 10 users to communicate pairwise using symmetric-key cryptography and using public-key cryptography respectively?
\optsshort{45 and 10}{100 and 20}{90 and 10}{45 and 20}
\end{question}

\begin{question}{H}
To create a digital signature the sender signs the message digest with
\optsshort{the sender's private key}{the sender's public key}{the receiver's private key}{the receiver's public key}
\end{question}

\begin{question}{H}
To ensure confidentiality in public-key cryptography, the sender encrypts the message with
\optsshort{the sender's private key}{the receiver's public key}{the sender's public key}{the receiver's private key}
\end{question}

\begin{question}{H}
A digital signature provides which of the following services?\\
A. Message authentication  B. Message integrity  C. Non-repudiation  D. Confidentiality\\
Choose the correct answer from the options given below:
\optsshort{B, C and D only}{A, B and C only}{A, B, C and D}{A and D only}
\end{question}

\begin{question}{H}
Match List I with List II.\\
List I (Parameter of DES): A. Block size  B. Effective key size  C. Number of rounds  D. Round key size\\
List II (Value): I. 16  II. 48 bits  III. 64 bits  IV. 56 bits
\optsshort{A-III, B-IV, C-II, D-I}{A-III, B-I, C-IV, D-II}{A-IV, B-III, C-I, D-II}{A-III, B-IV, C-I, D-II}
\end{question}

\begin{question}{H}
The number of rounds in AES for key sizes of 128, 192 and 256 bits are respectively
\optsshort{10, 12 and 14}{12, 14 and 16}{16, 16 and 16}{8, 10 and 12}
\end{question}

\begin{question}{H}
Which of the following is NOT a transformation in a round of AES?
\optsshort{ShiftRows}{SubBytes}{MixColumns}{Initial permutation}
\end{question}

\begin{question}{M}
Given below are two statements.\\
Statement I: DES is based on the Feistel cipher structure.\\
Statement II: AES is also based on the Feistel cipher structure.
\optslong{Both Statement I and Statement II are true}{Both Statement I and Statement II are false}{Statement I is true but Statement II is false}{Statement I is false but Statement II is true}
\end{question}

\begin{question}{H}
Which of the following are asymmetric-key algorithms?\\
A. RSA  B. AES  C. ElGamal  D. Blowfish  E. Elliptic curve cryptography\\
Choose the correct answer from the options given below:
\optsshort{A, B and C only}{A, C, D and E only}{A, C and E only}{B and D only}
\end{question}

\begin{question}{H}
The security of RSA and of Diffie-Hellman rests respectively on the difficulty of
\optslong{sorting and searching}{the knapsack problem and factoring}{factoring large integers and computing discrete logarithms}{computing discrete logarithms and factoring large integers}
\end{question}

\begin{question}{H}
The basic Diffie-Hellman key exchange is vulnerable to
\optsshort{a SQL injection attack}{a brute-force attack on DES}{a man-in-the-middle attack}{a birthday attack}
\end{question}

\begin{question}{H}
Using the Caesar cipher with a key of 3, the plaintext HELLO is encrypted as
\optsshort{IFMMP}{KHOOS}{KHOOR}{EBIIL}
\end{question}

\begin{question}{M}
Match List I with List II.\\
List I (Cipher): A. Caesar cipher  B. Rail fence cipher  C. Vigen\`ere cipher  D. Playfair cipher\\
List II (Type): I. Polyalphabetic substitution  II. Digram substitution using a 5 $\times$ 5 matrix  III. Monoalphabetic substitution  IV. Transposition
\optsshort{A-III, B-IV, C-II, D-I}{A-III, B-IV, C-I, D-II}{A-IV, B-III, C-I, D-II}{A-I, B-IV, C-III, D-II}
\end{question}

\begin{question}{M}
Which cipher is theoretically unbreakable when the key is random, as long as the message and used only once?
\optsshort{Playfair cipher}{Vigen\`ere cipher}{DES}{One-time pad}
\end{question}

\begin{question}{H}
Match List I with List II.\\
List I (Hash function): A. MD5  B. SHA-1  C. SHA-256  D. SHA-512\\
List II (Digest size): I. 160 bits  II. 512 bits  III. 128 bits  IV. 256 bits
\optsshort{A-III, B-I, C-II, D-IV}{A-III, B-IV, C-I, D-II}{A-I, B-III, C-IV, D-II}{A-III, B-I, C-IV, D-II}
\end{question}

\begin{question}{M}
A hash function produces an n-bit digest. The effort needed to find a collision using a birthday attack is of the order of
\optsshort{2\textsuperscript{n}}{2\textsuperscript{n/2}}{n\textsuperscript{2}}{2n}
\end{question}

\begin{question}{M}
Which property of a cryptographic hash function states that it is infeasible to find two different messages with the same digest?
\optsshort{Collision resistance}{Preimage resistance}{Second preimage resistance}{Avalanche effect}
\end{question}

\begin{question}{H}
Match List I with List II.\\
List I (Security protocol): A. IPSec  B. SSL/TLS  C. PGP  D. WPA2\\
List II (Where it is applied): I. E-mail security  II. Wireless LAN security  III. Network layer  IV. Between the transport and application layers
\optsshort{A-III, B-IV, C-I, D-II}{A-III, B-I, C-IV, D-II}{A-IV, B-III, C-I, D-II}{A-III, B-IV, C-II, D-I}
\end{question}

\begin{question}{H}
Given below are two statements.\\
Statement I: The Authentication Header (AH) protocol of IPSec provides source authentication and data integrity but not confidentiality.\\
Statement II: The Encapsulating Security Payload (ESP) protocol of IPSec provides confidentiality.
\optslong{Both Statement I and Statement II are true}{Both Statement I and Statement II are false}{Statement I is true but Statement II is false}{Statement I is false but Statement II is true}
\end{question}

\begin{question}{M}
In which IPSec mode is the entire original IP packet, including its header, protected and encapsulated in a new IP packet?
\optsshort{Tunnel mode}{Aggressive mode}{Main mode}{Transport mode}
\end{question}

\begin{question}{H}
A virtual private network (VPN) is
\optslong{a physically separate leased-line network}{a private network built over a public network using tunnelling and encryption}{a network without any firewall}{a LAN that uses only private IP addresses}
\end{question}

\begin{question}{H}
Match List I with List II.\\
List I (Malware): A. Virus  B. Worm  C. Trojan horse  D. Ransomware\\
List II (Characteristic): I. Encrypts the victim's files and demands payment  II. Appears to be a useful program but hides a malicious function  III. Attaches itself to a host program and needs it to spread  IV. Self-replicating, spreads across the network without a host program
\optsshort{A-III, B-IV, C-I, D-II}{A-III, B-II, C-IV, D-I}{A-III, B-IV, C-II, D-I}{A-IV, B-III, C-II, D-I}
\end{question}

\begin{question}{M}
A piece of malicious code that lies dormant until a specific date or condition occurs is called a
\optsshort{keylogger}{logic bomb}{rootkit}{worm}
\end{question}

\begin{question}{M}
Which type of malware hides its presence by modifying the operating system so that it retains privileged access?
\optsshort{Adware}{Rootkit}{Spyware}{Worm}
\end{question}

\begin{question}{H}
Which of the following are passive attacks?\\
A. Traffic analysis  B. Masquerade  C. Release of message contents (eavesdropping)  D. Replay  E. Denial of service\\
Choose the correct answer from the options given below:
\optsshort{C and D only}{B, D and E only}{A, B and C only}{A and C only}
\end{question}

\begin{question}{H}
Match List I with List II.\\
List I (Attack): A. Eavesdropping  B. Modification of messages  C. Denial of service  D. Masquerade\\
List II (Security goal threatened): I. Availability  II. Authentication  III. Confidentiality  IV. Integrity
\optsshort{A-III, B-IV, C-II, D-I}{A-IV, B-III, C-I, D-II}{A-III, B-I, C-IV, D-II}{A-III, B-IV, C-I, D-II}
\end{question}

\begin{question}{H}
Match List I with List II.\\
List I (Firewall type): A. Packet-filtering firewall  B. Stateful inspection firewall  C. Application-level gateway  D. Circuit-level gateway\\
List II (Characteristic): I. Acts as a proxy and inspects application data  II. Relays TCP connections without examining their contents  III. Filters each packet on IP addresses and port numbers  IV. Keeps track of the state of active connections
\optsshort{A-III, B-IV, C-II, D-I}{A-III, B-IV, C-I, D-II}{A-IV, B-III, C-I, D-II}{A-I, B-IV, C-III, D-II}
\end{question}

\begin{question}{M}
A packet-filtering firewall typically makes its decisions using
\optslong{digital signatures}{the contents of e-mail messages}{source and destination IP addresses, port numbers and protocol type}{user passwords}
\end{question}

\begin{question}{M}
A perimeter network that hosts an organisation's public servers, separated from both the Internet and the internal network, is called a
\optsshort{demilitarised zone (DMZ)}{virtual LAN}{honeypot}{extranet}
\end{question}

\begin{question}{M}
A decoy system designed to attract attackers and study their behaviour is called a
\optsshort{bastion host}{firewall}{proxy}{honeypot}
\end{question}

\begin{question}{H}
Steganography is the technique of
\optslong{converting plaintext to ciphertext using a key}{hiding the very existence of a message inside another medium such as an image}{compressing data before transmission}{verifying the identity of the sender}
\end{question}

\begin{question}{M}
A common technique for image steganography is
\optslong{increasing the image resolution}{encrypting the image with RSA}{modifying the least significant bits of pixel values}{changing the image file name}
\end{question}

\begin{question}{H}
Given below are two statements: one is labelled as Assertion A and the other is labelled as Reason R.\\
Assertion A: In practice public-key algorithms are used to exchange a session key rather than to encrypt bulk data.\\
Reason R: Public-key algorithms are computationally much slower than symmetric-key algorithms.
\optslong{Both A and R are true and R is the correct explanation of A}{Both A and R are true but R is NOT the correct explanation of A}{A is true but R is false}{A is false but R is true}
\end{question}

\begin{question}{M}
Which of the following binds a public key to the identity of its owner and is signed by a certification authority?
\optsshort{Message digest}{Session key}{X.509 digital certificate}{Nonce}
\end{question}

\begin{question}{M}
Kerberos is
\optslong{a hashing algorithm}{a firewall architecture}{a public-key encryption algorithm}{an authentication protocol based on symmetric keys, tickets and a key distribution centre}
\end{question}

\begin{question}{M}
A nonce (a number used only once) or a timestamp in an authentication protocol protects against
\optsshort{traffic analysis}{dictionary attacks}{replay attacks}{brute-force attacks}
\end{question}

\begin{question}{M}
The value of Euler's totient function $\varphi$(35) is
\optsshort{24}{35}{12}{34}
\end{question}

\begin{question}{M}
Which block cipher mode of operation encrypts identical plaintext blocks into identical ciphertext blocks?
\optsshort{CBC}{ECB}{CFB}{CTR}
\end{question}

\begin{question}{M}
Triple DES with three independent keys has a key length of
\optsshort{56 bits}{112 bits}{128 bits}{168 bits}
\end{question}

\begin{question}{M}
Which of the following is a stream cipher?
\optsshort{IDEA}{AES}{RC4}{DES}
\end{question}

\begin{question}{M}
According to Shannon, the two properties that a good cipher should have are
\optsshort{authentication and integrity}{speed and simplicity}{substitution and compression}{confusion and diffusion}
\end{question}

\begin{question}{M}
Kerckhoffs's principle states that the security of a cryptosystem should depend on
\optsshort{the secrecy of the algorithm}{the length of the plaintext}{the speed of encryption}{the secrecy of the key only}
\end{question}

\begin{question}{M}
An attack in which fraudulent e-mails or websites trick users into revealing passwords or card details is called
\optsshort{war driving}{sniffing}{spoofing of MAC addresses}{phishing}
\end{question}

\begin{question}{M}
Which of the following combines a secret key with a hash function to provide message authentication?
\optsshort{Base64}{HMAC}{MD5 alone}{CRC}
\end{question}

\begin{question}{M}
In a distributed denial-of-service attack
\optslong{one attacker guesses passwords of many users}{data is silently copied from a server}{many compromised machines flood a target to make it unavailable}{a certificate is forged}
\end{question}

\begin{question}{M}
Given below are two statements.\\
Statement I: An intrusion detection system (IDS) monitors traffic and raises alerts on suspicious activity.\\
Statement II: An intrusion prevention system (IPS) is placed in-line and can block malicious traffic.
\optslong{Both Statement I and Statement II are true}{Both Statement I and Statement II are false}{Statement I is true but Statement II is false}{Statement I is false but Statement II is true}
\end{question}

\begin{question}{M}
Which wireless security protocol uses AES-based CCMP for encryption?
\optsshort{WEP}{WPA (TKIP)}{WPA2}{None of these}
\end{question}

\begin{question}{L}
An attack that tries every possible key until the correct one is found is a
\optsshort{side-channel attack}{chosen-plaintext attack}{replay attack}{brute-force attack}
\end{question}

\begin{question}{M}
In which cryptanalytic attack does the attacker have access to pairs of plaintext and the corresponding ciphertext, but cannot choose them?
\optsshort{Ciphertext-only attack}{Chosen-ciphertext attack}{Known-plaintext attack}{Chosen-plaintext attack}
\end{question}

\begin{answers}
\ans{1}{2}{$\varphi$(n) = 2 $\times$ 10 = 20; 7 $\times$ 3 = 21 $\equiv$ 1 (mod 20), so d = 3.}
\ans{2}{1}{C = 2\textsuperscript{7} mod 33 = 128 mod 33 = 29.}
\ans{3}{2}{$\varphi$(n) = 60; 13 $\times$ 37 = 481 = 8 $\times$ 60 + 1.}
\ans{4}{1}{$\varphi$(n) = 40; 3 $\times$ 27 = 81 $\equiv$ 1 (mod 40).}
\ans{5}{4}{5\textsuperscript{6} mod 23: 5\textsuperscript{2} = 2, 5\textsuperscript{4} = 4, so 5\textsuperscript{6} = 8.}
\ans{6}{3}{B's public key is 5\textsuperscript{15} mod 23 = 19; K = 19\textsuperscript{6} mod 23 = 2.}
\ans{7}{4}{Symmetric: n(n $-$ 1)/2 = 45 shared keys; asymmetric: one key pair per user = 20 keys.}
\ans{8}{1}{Anyone can verify it with the sender's public key.}
\ans{9}{2}{Only the receiver's private key can decrypt it.}
\ans{10}{2}{A signature alone does not hide the message; confidentiality needs encryption.}
\ans{11}{4}{DES: 64-bit block, 56-bit key, 16 Feistel rounds, 48-bit round keys.}
\ans{12}{1}{AES always has a 128-bit block; the rounds depend on the key size.}
\ans{13}{4}{AES rounds use SubBytes, ShiftRows, MixColumns and AddRoundKey; the initial permutation belongs to DES.}
\ans{14}{3}{AES (Rijndael) is a substitution-permutation network, not a Feistel cipher.}
\ans{15}{3}{AES and Blowfish are symmetric block ciphers.}
\ans{16}{3}{RSA: integer factorisation; Diffie-Hellman and ElGamal: discrete logarithm.}
\ans{17}{3}{It does not authenticate the two parties, so an attacker can sit between them.}
\ans{18}{3}{Each letter is shifted forward by three positions.}
\ans{19}{2}{Caesar is monoalphabetic, rail fence rearranges letters, Vigen\`ere is polyalphabetic, Playfair encrypts letter pairs.}
\ans{20}{4}{The one-time pad (Vernam cipher) offers perfect secrecy.}
\ans{21}{4}{MD5 128, SHA-1 160, SHA-256 256, SHA-512 512 bits.}
\ans{22}{2}{The birthday paradox halves the exponent for collision search.}
\ans{23}{1}{Collision resistance is also called strong collision resistance.}
\ans{24}{1}{IPSec secures IP, TLS secures transport connections, PGP secures mail, WPA2 secures Wi-Fi.}
\ans{25}{1}{ESP can provide encryption as well as authentication; AH cannot encrypt.}
\ans{26}{1}{Tunnel mode is used between gateways, e.g. in VPNs; transport mode protects only the payload.}
\ans{27}{2}{VPNs give privacy over the Internet using protocols such as IPSec, L2TP and PPTP.}
\ans{28}{3}{Standard definitions of the four malware types.}
\ans{29}{2}{A logic bomb is triggered by an event such as a date or the deletion of a record.}
\ans{30}{2}{A rootkit conceals itself and other malware at the kernel or system level.}
\ans{31}{4}{Passive attacks only observe; masquerade, replay, modification and denial of service are active.}
\ans{32}{4}{The CIA triad plus authentication.}
\ans{33}{2}{Standard classification of firewalls.}
\ans{34}{3}{It examines network- and transport-layer headers only.}
\ans{35}{1}{The DMZ limits the damage if a public-facing server is compromised.}
\ans{36}{4}{A honeypot diverts attackers from real systems and gathers information about them.}
\ans{37}{2}{Cryptography hides the meaning of a message; steganography hides its existence.}
\ans{38}{3}{LSB substitution changes the image imperceptibly.}
\ans{39}{1}{Hybrid cryptosystems combine the convenience of public-key distribution with the speed of symmetric encryption.}
\ans{40}{3}{X.509 certificates are the basis of public-key infrastructure.}
\ans{41}{4}{Kerberos uses an authentication server and a ticket-granting server.}
\ans{42}{3}{An old captured message cannot be reused because its nonce is no longer fresh.}
\ans{43}{1}{35 = 5 $\times$ 7, so $\varphi$ = 4 $\times$ 6 = 24.}
\ans{44}{2}{Electronic codebook mode encrypts each block independently, so patterns leak.}
\ans{45}{4}{3 $\times$ 56 = 168 bits (112 bits with two keys).}
\ans{46}{3}{RC4 generates a key stream that is XORed with the plaintext; the others are block ciphers.}
\ans{47}{4}{Confusion hides the key-ciphertext relationship; diffusion spreads plaintext statistics.}
\ans{48}{4}{The algorithm is assumed to be public.}
\ans{49}{4}{Phishing is a social-engineering attack.}
\ans{50}{2}{An HMAC proves both integrity and that the sender knows the shared key.}
\ans{51}{3}{A botnet of zombies overwhelms the victim's resources.}
\ans{52}{1}{Detection versus active prevention.}
\ans{53}{3}{WEP and TKIP use RC4; WPA2 mandates AES-CCMP.}
\ans{54}{4}{For DES this needs at most 2\textsuperscript{56} trials.}
\ans{55}{3}{In a chosen-plaintext attack the analyst additionally selects the plaintexts.}
\end{answers}

\section{Mobile Technology (cellular systems, GSM/GPRS/SMS, wireless LANs, MANETs, satellites, geolocation)}

\begin{question}{H}
In a cellular system with hexagonal cells, the cluster size N must satisfy N = i\textsuperscript{2} + ij + j\textsuperscript{2} for non-negative integers i and j. Which of the following is NOT a valid cluster size?
\optsshort{5}{3}{4}{7}
\end{question}

\begin{question}{H}
A cellular system has 490 duplex channels and uses a cluster size of 7. The number of channels available per cell is
\optsshort{490}{98}{35}{70}
\end{question}

\begin{question}{H}
For a cluster size of N = 7 and cell radius R, the co-channel reuse distance D is approximately
\optsshort{2.65 R}{7 R}{4.58 R}{3.46 R}
\end{question}

\begin{question}{H}
Given below are two statements.\\
Statement I: Decreasing the cluster size N increases the capacity of a cellular system.\\
Statement II: Decreasing the cluster size N reduces co-channel interference.
\optslong{Both Statement I and Statement II are true}{Both Statement I and Statement II are false}{Statement I is true but Statement II is false}{Statement I is false but Statement II is true}
\end{question}

\begin{question}{H}
The process of transferring an ongoing call from one base station to another as the mobile moves is called
\optsshort{roaming}{registration}{paging}{handoff (handover)}
\end{question}

\begin{question}{H}
Given below are two statements.\\
Statement I: In a hard handoff the mobile breaks the connection with the old base station before connecting to the new one.\\
Statement II: Soft handoff, in which the mobile communicates with two base stations simultaneously, is a feature of CDMA systems.
\optslong{Both Statement I and Statement II are true}{Both Statement I and Statement II are false}{Statement I is true but Statement II is false}{Statement I is false but Statement II is true}
\end{question}

\begin{question}{M}
Which techniques are used to increase the capacity of a cellular system?\\
A. Cell splitting  B. Sectoring  C. Increasing the cell radius  D. Microcell zone concept\\
Choose the correct answer from the options given below:
\optsshort{B, C and D only}{A and C only}{A, B, C and D}{A, B and D only}
\end{question}

\begin{question}{H}
Match List I with List II.\\
List I (Generation): A. 1G  B. 2G  C. 3G  D. 4G\\
List II (Representative technology): I. LTE  II. UMTS / WCDMA  III. AMPS  IV. GSM
\optsshort{A-III, B-II, C-IV, D-I}{A-III, B-IV, C-I, D-II}{A-IV, B-III, C-II, D-I}{A-III, B-IV, C-II, D-I}
\end{question}

\begin{question}{H}
GSM uses which multiple access technique?
\optslong{CDMA}{FDMA only}{A combination of FDMA and TDMA}{OFDMA}
\end{question}

\begin{question}{M}
In GSM, the carrier spacing and the number of time slots per TDMA frame are respectively
\optsshort{1.25 MHz and 64}{30 kHz and 6}{200 kHz and 8}{25 kHz and 3}
\end{question}

\begin{question}{M}
The modulation technique used in GSM is
\optsshort{QPSK}{BFSK}{GMSK}{16-QAM}
\end{question}

\begin{question}{H}
Match List I with List II.\\
List I (GSM entity): A. HLR  B. VLR  C. EIR  D. AuC\\
List II (Function): I. Stores authentication and ciphering keys  II. Database of valid and stolen mobile equipment (IMEI)  III. Permanent database of subscribers of the home network  IV. Temporary database of subscribers currently in an MSC area
\optsshort{A-IV, B-III, C-II, D-I}{A-III, B-IV, C-II, D-I}{A-III, B-IV, C-I, D-II}{A-III, B-II, C-IV, D-I}
\end{question}

\begin{question}{M}
Given below are two statements.\\
Statement I: The IMSI identifies the subscriber and is stored in the SIM.\\
Statement II: The IMEI identifies the mobile equipment.
\optslong{Both Statement I and Statement II are true}{Both Statement I and Statement II are false}{Statement I is true but Statement II is false}{Statement I is false but Statement II is true}
\end{question}

\begin{question}{M}
In GSM, the algorithms A3, A5 and A8 are used respectively for
\optslong{encryption, authentication and key generation}{authentication, encryption and cipher key generation}{authentication, key generation and encryption}{key generation, authentication and encryption}
\end{question}

\begin{question}{M}
Arrange the GSM network elements in the order in which a call from a mobile station passes through them.\\
A. BSC  B. MSC  C. BTS
\optsshort{C, B, A}{B, A, C}{A, C, B}{C, A, B}
\end{question}

\begin{question}{H}
GPRS adds packet-switched data to GSM by introducing which two new network nodes?
\optsshort{HLR and VLR}{MSC and GMSC}{SGSN and GGSN}{BSC and BTS}
\end{question}

\begin{question}{M}
GPRS and EDGE are commonly referred to respectively as
\optsshort{2.5G and 2.75G}{1G and 2G}{2G and 3G}{3G and 4G}
\end{question}

\begin{question}{H}
The maximum length of a single SMS message using the default 7-bit GSM alphabet is
\optsshort{256 characters}{140 characters}{160 characters}{70 characters}
\end{question}

\begin{question}{M}
SMS messages are delivered using
\optslong{GPRS tunnels only}{Wi-Fi access points}{signalling channels and a store-and-forward SMS centre}{a dedicated traffic channel and circuit switching}
\end{question}

\begin{question}{H}
Given below are two statements: one is labelled as Assertion A and the other is labelled as Reason R.\\
Assertion A: CDMA systems require tight power control.\\
Reason R: Without it, a transmitter near the base station would drown out the signals of distant mobiles (the near-far problem).
\optslong{Both A and R are true and R is the correct explanation of A}{Both A and R are true but R is NOT the correct explanation of A}{A is true but R is false}{A is false but R is true}
\end{question}

\begin{question}{M}
The multiple access technique used in the downlink of 4G LTE is
\optsshort{CDMA}{OFDMA}{FDMA}{TDMA}
\end{question}

\begin{question}{H}
Match List I with List II.\\
List I (Standard): A. IEEE 802.11a  B. IEEE 802.11b  C. IEEE 802.11g  D. IEEE 802.11n\\
List II (Band and maximum data rate): I. 2.4 GHz, 54 Mbps  II. 2.4/5 GHz, up to 600 Mbps using MIMO  III. 5 GHz, 54 Mbps  IV. 2.4 GHz, 11 Mbps
\optsshort{A-III, B-IV, C-II, D-I}{A-III, B-IV, C-I, D-II}{A-IV, B-III, C-I, D-II}{A-I, B-IV, C-III, D-II}
\end{question}

\begin{question}{M}
An IEEE 802.11 network that operates without an access point is called
\optslong{an independent BSS (ad hoc network)}{an extended service set}{a distribution system}{an infrastructure BSS}
\end{question}

\begin{question}{M}
Two or more basic service sets interconnected by a distribution system form
\optsshort{an IBSS}{a scatternet}{a piconet}{an ESS}
\end{question}

\begin{question}{M}
In IEEE 802.11, the distributed coordination function (DCF) uses
\optsshort{polling}{CSMA/CA}{token passing}{CSMA/CD}
\end{question}

\begin{question}{H}
Station A and station C are both within range of station B but cannot hear each other. If A and C transmit to B at the same time, the collision at B is an example of the
\optsshort{hidden terminal problem}{capture effect}{near-far problem}{exposed terminal problem}
\end{question}

\begin{question}{H}
A Bluetooth piconet can have
\optslong{any number of masters}{2 masters and up to 6 active slaves}{1 master and up to 8 active slaves}{1 master and up to 7 active slaves}
\end{question}

\begin{question}{M}
Bluetooth operates in the 2.4 GHz ISM band and uses
\optslong{OFDM}{DSSS}{frequency hopping spread spectrum at 1600 hops per second}{infrared}
\end{question}

\begin{question}{M}
Two or more piconets linked by a device that participates in both form a
\optsshort{scatternet}{basic service set}{cluster}{mesh cell}
\end{question}

\begin{question}{M}
Match List I with List II.\\
List I (Technology): A. Bluetooth  B. ZigBee  C. Wi-Fi  D. WiMAX\\
List II (IEEE standard): I. 802.16  II. 802.11  III. 802.15.1  IV. 802.15.4
\optsshort{A-IV, B-III, C-II, D-I}{A-III, B-IV, C-I, D-II}{A-III, B-II, C-IV, D-I}{A-III, B-IV, C-II, D-I}
\end{question}

\begin{question}{H}
In Mobile IP, the temporary address that a mobile node acquires in a foreign network is called the
\optsshort{link-local address}{anycast address}{home address}{care-of address}
\end{question}

\begin{question}{H}
In Mobile IP, packets addressed to the mobile node's home address are intercepted by the
\optslong{correspondent node}{foreign agent and tunnelled to the home agent}{home agent and tunnelled to the care-of address}{DNS server}
\end{question}

\begin{question}{M}
The inefficiency in Mobile IP in which packets from the correspondent node travel via the home agent even when the mobile node is nearby is called
\optsshort{count to infinity}{triangle routing}{tunnel collapse}{ping-pong handoff}
\end{question}

\begin{question}{H}
A mobile ad hoc network (MANET) is characterised by
\optslong{a fixed infrastructure of base stations}{self-configuring mobile nodes that act as both hosts and routers without fixed infrastructure}{a static topology}{wired backbone links between nodes}
\end{question}

\begin{question}{H}
Match List I with List II.\\
List I (MANET routing protocol): A. DSDV  B. AODV  C. ZRP\\
List II (Category): I. Hybrid  II. Proactive (table-driven)  III. Reactive (on-demand)
\optsshort{A-I, B-III, C-II}{A-II, B-III, C-I}{A-II, B-I, C-III}{A-III, B-II, C-I}
\end{question}

\begin{question}{M}
In the Dynamic Source Routing (DSR) protocol
\optslong{routes are computed by a central server}{every node maintains a route to every other node at all times}{link state packets are flooded periodically}{the complete route is carried in the header of each data packet}
\end{question}

\begin{question}{M}
In AODV, a route to a destination is discovered by
\optslong{flooding the data packets}{broadcasting a route request (RREQ) and receiving a route reply (RREP)}{consulting the home agent}{periodic exchange of full routing tables}
\end{question}

\begin{question}{H}
Match List I with List II.\\
List I (Satellite orbit): A. LEO  B. MEO  C. GEO\\
List II (Approximate altitude): I. 35,786 km  II. 500 to 2,000 km  III. 5,000 to 20,000 km
\optsshort{A-II, B-III, C-I}{A-I, B-III, C-II}{A-II, B-I, C-III}{A-III, B-II, C-I}
\end{question}

\begin{question}{M}
The minimum number of geostationary satellites needed to cover almost the whole Earth (excluding the polar regions) is
\optsshort{24}{6}{3}{2}
\end{question}

\begin{question}{M}
The round-trip propagation delay (Earth station to satellite and back to Earth) of a geostationary satellite link is approximately
\optsshort{25 ms}{2.5 ms}{2.5 s}{250 ms}
\end{question}

\begin{question}{M}
Given below are two statements.\\
Statement I: LEO satellites give lower propagation delay than GEO satellites.\\
Statement II: A LEO system needs fewer satellites than a GEO system for global coverage.
\optslong{Both Statement I and Statement II are true}{Both Statement I and Statement II are false}{Statement I is true but Statement II is false}{Statement I is false but Statement II is true}
\end{question}

\begin{question}{M}
The equipment on a satellite that receives the uplink signal, amplifies it, changes its frequency and retransmits it on the downlink is the
\optsshort{multiplexer}{duplexer}{transponder}{footprint}
\end{question}

\begin{question}{H}
The Global Positioning System nominally consists of
\optslong{3 geostationary satellites}{24 satellites in 6 orbital planes}{12 satellites in 3 orbital planes}{66 satellites in 6 orbital planes}
\end{question}

\begin{question}{H}
The minimum number of GPS satellites that a receiver needs in order to compute its three-dimensional position and correct its clock is
\optsshort{6}{2}{4}{3}
\end{question}

\begin{question}{M}
GPS determines the position of a receiver by
\optslong{measuring signal strength of cell towers}{trilateration using distances derived from signal travel times}{Doppler shift only}{triangulation using angles of arrival}
\end{question}

\begin{question}{M}
Which of the following is a network-based wireless geolocation technique used in cellular systems?
\optslong{Frequency reuse}{Slow start}{Time difference of arrival (TDOA)}{Bit stuffing}
\end{question}

\begin{question}{M}
Which of the following has the shortest range?
\optsshort{WiMAX}{Wi-Fi}{Bluetooth}{NFC}
\end{question}

\begin{question}{M}
WAP (Wireless Application Protocol) uses which markup language for content on small-screen devices?
\optsshort{HTML}{SGML}{WML}{XAML}
\end{question}

\begin{question}{M}
In first-generation analog cellular systems such as AMPS, the multiple access technique used was
\optsshort{TDMA}{CDMA}{FDMA}{OFDMA}
\end{question}

\begin{question}{M}
Why are cells in a cellular network modelled as hexagons?
\optslong{A hexagon approximates a circle and tiles a plane without gaps or overlaps}{Antennas radiate in a hexagonal pattern}{Hexagons need the least transmission power}{It is required by the GSM standard}
\end{question}

\begin{question}{M}
Interference between cells that use the same set of frequencies is called
\optslong{adjacent-channel interference}{crosstalk}{intersymbol interference}{co-channel interference}
\end{question}

\begin{question}{M}
Which of the following is a feature of 5G networks?\\
A. Millimetre-wave bands  B. Massive MIMO  C. Network slicing  D. Circuit-switched voice core\\
Choose the correct answer from the options given below:
\optsshort{B, C and D only}{A, B and C only}{A and D only}{A, B, C and D}
\end{question}

\begin{question}{L}
The 2.4 GHz band used by Wi-Fi and Bluetooth is part of the
\optslong{licensed cellular band}{C band}{Ku band}{ISM (industrial, scientific and medical) band}
\end{question}

\begin{question}{M}
In IEEE 802.11, a station that hears an RTS or CTS frame sets its \underline{\hspace{3em}} and defers access for the indicated duration.
\optslong{backoff counter to zero}{beacon interval}{network allocation vector (NAV)}{contention window}
\end{question}

\begin{answers}
\ans{1}{1}{Valid values are 1, 3, 4, 7, 9, 12, 13 ...; 5 cannot be written in this form.}
\ans{2}{4}{490 / 7 = 70 channels per cell.}
\ans{3}{3}{D = R $\times$ sqrt(3N) = R $\times$ sqrt(21) $\approx$ 4.58 R.}
\ans{4}{3}{A smaller N gives more channels per cell but brings co-channel cells closer, increasing interference.}
\ans{5}{4}{Roaming refers to service in a network other than the home network.}
\ans{6}{1}{Hard handoff is break-before-make; soft handoff is make-before-break.}
\ans{7}{4}{Larger cells reduce capacity per unit area; splitting, sectoring and microzones increase it.}
\ans{8}{4}{1G analog AMPS, 2G digital GSM, 3G UMTS, 4G LTE.}
\ans{9}{3}{Each 200 kHz carrier is divided into 8 time slots.}
\ans{10}{3}{GSM carriers are 200 kHz wide with 8 full-rate slots.}
\ans{11}{3}{Gaussian minimum shift keying gives a compact spectrum with constant envelope.}
\ans{12}{2}{Home/visitor location registers, equipment identity register, authentication centre.}
\ans{13}{1}{IMSI is tied to the subscription, IMEI to the handset.}
\ans{14}{2}{A3 authenticates, A8 derives the session key, A5 encrypts the radio link.}
\ans{15}{4}{Mobile station, base transceiver station, base station controller, mobile switching centre.}
\ans{16}{3}{The serving and gateway GPRS support nodes route packets and connect to external data networks.}
\ans{17}{1}{GPRS is 2.5G; EDGE, with 8-PSK modulation, is 2.75G.}
\ans{18}{3}{140 bytes of payload hold 160 seven-bit characters (70 in 16-bit UCS-2).}
\ans{19}{3}{The SMSC stores the message and forwards it when the recipient is reachable.}
\ans{20}{1}{All users share the same frequency, so received powers must be equalised.}
\ans{21}{2}{LTE uses OFDMA on the downlink and SC-FDMA on the uplink.}
\ans{22}{2}{802.11a: 5 GHz/54; b: 2.4 GHz/11; g: 2.4 GHz/54; n: MIMO up to 600 Mbps.}
\ans{23}{1}{In an IBSS the stations communicate directly with one another.}
\ans{24}{4}{An extended service set lets stations roam between access points.}
\ans{25}{2}{DCF is contention-based CSMA/CA; the optional PCF uses polling by the access point.}
\ans{26}{1}{Each sender is hidden from the other; RTS/CTS mitigates it.}
\ans{27}{4}{3-bit active member addresses allow 7 active slaves; more can be parked.}
\ans{28}{3}{Classic Bluetooth hops over 79 channels of 1 MHz.}
\ans{29}{1}{A bridge node may be a slave in one piconet and a master or slave in another.}
\ans{30}{4}{Bluetooth 802.15.1, ZigBee 802.15.4, Wi-Fi 802.11, WiMAX 802.16.}
\ans{31}{4}{The home agent tunnels packets to the care-of address.}
\ans{32}{3}{The home agent encapsulates the packets (IP-in-IP) and forwards them to the foreign network.}
\ans{33}{2}{Route optimisation lets the correspondent node send directly to the care-of address.}
\ans{34}{2}{Topology changes dynamically and routing is multi-hop.}
\ans{35}{2}{DSDV and OLSR are proactive, AODV and DSR are reactive, ZRP is hybrid.}
\ans{36}{4}{DSR is an on-demand source-routing protocol; AODV instead keeps next-hop entries at the nodes.}
\ans{37}{2}{Route discovery is triggered only when a route is needed.}
\ans{38}{1}{Low, medium and geostationary Earth orbits.}
\ans{39}{3}{Three satellites spaced 120 degrees apart suffice.}
\ans{40}{4}{About 36,000 km each way at 3 $\times$ 10\textsuperscript{8} m/s gives 240 to 270 ms.}
\ans{41}{3}{LEO footprints are small and the satellites move quickly, so many satellites and frequent handoffs are needed.}
\ans{42}{3}{The uplink frequency is higher than the downlink frequency to avoid interference.}
\ans{43}{2}{GPS satellites are in MEO at about 20,200 km.}
\ans{44}{3}{Three ranges give position (trilateration); the fourth resolves the receiver clock error.}
\ans{45}{2}{Each range is the travel time multiplied by the speed of light.}
\ans{46}{3}{TDOA, angle of arrival and cell-ID are common network-based positioning methods.}
\ans{47}{4}{Near-field communication works over about 4 cm or less.}
\ans{48}{3}{Wireless Markup Language organises content into decks and cards.}
\ans{49}{3}{Each call was given its own 30 kHz frequency channel.}
\ans{50}{1}{Among shapes that tessellate, the hexagon covers the largest area for a given radius.}
\ans{51}{4}{It is controlled by the reuse distance, not by raising transmit power.}
\ans{52}{2}{4G and 5G cores are all-IP and packet-switched.}
\ans{53}{4}{ISM bands are unlicensed.}
\ans{54}{3}{The NAV implements virtual carrier sensing.}
\end{answers}

\section{Cloud Computing and IoT (service and deployment models, virtualization, storage, SLA, IoT basics)}

\begin{question}{H}
Match List I with List II.\\
List I (Service model): A. SaaS  B. PaaS  C. IaaS\\
List II (Example): I. Amazon EC2  II. Gmail  III. Google App Engine
\optsshort{A-II, B-I, C-III}{A-III, B-II, C-I}{A-I, B-III, C-II}{A-II, B-III, C-I}
\end{question}

\begin{question}{H}
In which cloud service model does the consumer manage the operating system, runtime and applications, while the provider manages virtualization, servers, storage and networking?
\optsshort{FaaS}{PaaS}{IaaS}{SaaS}
\end{question}

\begin{question}{H}
A developer deploys application code to a cloud service that supplies the operating system, runtime, middleware and database, without managing the servers. This is
\optsshort{On-premises hosting}{SaaS}{PaaS}{IaaS}
\end{question}

\begin{question}{H}
Arrange the cloud service models in increasing order of the control the consumer has over the underlying resources.\\
A. IaaS  B. SaaS  C. PaaS
\optsshort{A, C, B}{B, A, C}{B, C, A}{C, B, A}
\end{question}

\begin{question}{H}
Which of the following are essential characteristics of cloud computing according to NIST?\\
A. On-demand self-service  B. Broad network access  C. Resource pooling  D. Rapid elasticity  E. Measured service\\
Choose the correct answer from the options given below:
\optsshort{A, B, C, D and E}{A, B and C only}{A, C and D only}{B, D and E only}
\end{question}

\begin{question}{H}
Match List I with List II.\\
List I (Deployment model): A. Public cloud  B. Private cloud  C. Hybrid cloud  D. Community cloud\\
List II (Description): I. Shared by several organisations with common concerns  II. Composition of two or more distinct clouds bound together  III. Open for use by the general public  IV. Operated solely for a single organisation
\optsshort{A-III, B-II, C-IV, D-I}{A-IV, B-III, C-II, D-I}{A-III, B-IV, C-I, D-II}{A-III, B-IV, C-II, D-I}
\end{question}

\begin{question}{M}
An organisation runs its normal workload in its private cloud and automatically uses a public cloud when demand spikes. This is called
\optsshort{multi-tenancy}{vendor lock-in}{cloud bursting}{live migration}
\end{question}

\begin{question}{H}
Given below are two statements.\\
Statement I: A Type 1 hypervisor runs directly on the host hardware.\\
Statement II: A Type 2 hypervisor runs as an application on top of a host operating system.
\optslong{Both Statement I and Statement II are true}{Both Statement I and Statement II are false}{Statement I is true but Statement II is false}{Statement I is false but Statement II is true}
\end{question}

\begin{question}{H}
Which of the following is a Type 2 (hosted) hypervisor?
\optsshort{Microsoft Hyper-V}{Xen}{VMware ESXi}{Oracle VirtualBox}
\end{question}

\begin{question}{H}
Given below are two statements.\\
Statement I: In full virtualization the guest operating system runs unmodified and is unaware that it is virtualized.\\
Statement II: In paravirtualization the guest operating system is modified to make hypercalls to the hypervisor.
\optslong{Both Statement I and Statement II are true}{Both Statement I and Statement II are false}{Statement I is true but Statement II is false}{Statement I is false but Statement II is true}
\end{question}

\begin{question}{M}
The software layer that creates and runs virtual machines by multiplexing physical hardware among them is the
\optslong{load balancer}{device driver}{hypervisor (virtual machine monitor)}{container registry}
\end{question}

\begin{question}{H}
Given below are two statements: one is labelled as Assertion A and the other is labelled as Reason R.\\
Assertion A: Containers are more lightweight than virtual machines.\\
Reason R: Containers share the host operating system kernel instead of each running a full guest operating system.
\optslong{Both A and R are true and R is the correct explanation of A}{Both A and R are true but R is NOT the correct explanation of A}{A is true but R is false}{A is false but R is true}
\end{question}

\begin{question}{M}
Moving a running virtual machine from one physical host to another without noticeable downtime is called
\optsshort{live migration}{snapshotting}{provisioning}{cloning}
\end{question}

\begin{question}{M}
The property by which a single instance of software serves several customers, with their data kept isolated, is called
\optsshort{multi-tenancy}{multiprogramming}{multithreading}{multihoming}
\end{question}

\begin{question}{H}
Given below are two statements.\\
Statement I: Vertical scaling (scaling up) adds more CPU or memory to an existing server.\\
Statement II: Horizontal scaling (scaling out) adds more servers or instances to the pool.
\optslong{Both Statement I and Statement II are true}{Both Statement I and Statement II are false}{Statement I is true but Statement II is false}{Statement I is false but Statement II is true}
\end{question}

\begin{question}{M}
The ability of a cloud system to acquire and release resources automatically as the workload changes is called
\optsshort{fault tolerance}{interoperability}{portability}{elasticity}
\end{question}

\begin{question}{H}
A formal contract between a cloud provider and a consumer that specifies guaranteed levels of availability, performance and penalties is a
\optslong{non-disclosure agreement}{service level agreement (SLA)}{service catalogue}{software licence}
\end{question}

\begin{question}{M}
An SLA guarantees 99.9\% availability. The maximum permitted downtime in a 30-day month is approximately
\optsshort{7.2 hours}{4.3 minutes}{43 minutes}{72 hours}
\end{question}

\begin{question}{H}
Match List I with List II.\\
List I (Storage type): A. Object storage  B. Block storage  C. File storage\\
List II (Example or characteristic): I. Shared hierarchical file system accessed through NFS or SMB  II. Amazon S3 --- data kept as objects with metadata in a flat namespace  III. Raw volumes attached to a virtual machine, e.g. Amazon EBS
\optsshort{A-II, B-I, C-III}{A-I, B-III, C-II}{A-III, B-II, C-I}{A-II, B-III, C-I}
\end{question}

\begin{question}{M}
A database offered as a managed cloud service, in which the provider handles installation, patching, backup and scaling, is called
\optsshort{DBaaS}{NaaS}{a data warehouse appliance}{DaaS (desktop as a service)}
\end{question}

\begin{question}{M}
According to the CAP theorem a distributed data store cannot simultaneously guarantee all three of
\optslong{consistency, availability and partition tolerance}{concurrency, atomicity and performance}{consistency, atomicity and persistence}{capacity, accuracy and privacy}
\end{question}

\begin{question}{M}
In the pay-as-you-go model of cloud computing
\optslong{billing is a fixed annual fee}{the customer buys the hardware}{operational expenditure is converted into capital expenditure}{capital expenditure is converted into operational expenditure}
\end{question}

\begin{question}{M}
Which of the following is a concern or disadvantage of cloud computing?
\optsshort{On-demand self-service}{Elasticity}{Reduced capital cost}{Vendor lock-in}
\end{question}

\begin{question}{M}
In serverless computing (Function as a Service)
\optslong{the customer must manage the virtual machines}{there are no servers involved}{applications run only on the client}{the provider manages server provisioning and the customer is billed per function execution}
\end{question}

\begin{question}{M}
Processing data close to where it is generated, instead of sending everything to a distant data centre, is called
\optsshort{edge (fog) computing}{grid computing}{utility computing}{batch computing}
\end{question}

\begin{question}{M}
Which of the following is an open-source platform for building private and public IaaS clouds?
\optsshort{Hadoop}{MongoDB}{OpenStack}{Kubernetes}
\end{question}

\begin{question}{M}
In cloud resource management, distributing incoming requests across multiple servers so that no single server is overloaded is called
\optsshort{load balancing}{orchestration of storage}{virtualization}{caching}
\end{question}

\begin{question}{M}
Which programming model, popularised by Google and implemented in Hadoop, processes large data sets in parallel across a cluster?
\optsshort{Publish-subscribe}{Two-phase commit}{MapReduce}{Remote procedure call}
\end{question}

\begin{question}{M}
Given below are two statements.\\
Statement I: Server virtualization allows several virtual servers to run on one physical server.\\
Statement II: Server consolidation through virtualization increases hardware utilisation and reduces energy consumption.
\optslong{Both Statement I and Statement II are true}{Both Statement I and Statement II are false}{Statement I is true but Statement II is false}{Statement I is false but Statement II is true}
\end{question}

\begin{question}{H}
The term ``Internet of Things'' was coined by
\optsshort{Tim Berners-Lee}{Kevin Ashton}{Vint Cerf}{John McCarthy}
\end{question}

\begin{question}{H}
Arrange the layers of the basic three-layer IoT architecture from bottom to top.\\
A. Network layer  B. Application layer  C. Perception layer
\optsshort{C, B, A}{C, A, B}{B, A, C}{A, C, B}
\end{question}

\begin{question}{H}
Given below are two statements.\\
Statement I: A sensor converts a physical quantity into an electrical signal.\\
Statement II: An actuator converts an electrical signal into a physical action.
\optslong{Both Statement I and Statement II are true}{Both Statement I and Statement II are false}{Statement I is true but Statement II is false}{Statement I is false but Statement II is true}
\end{question}

\begin{question}{H}
Which of the following statements about MQTT is correct?
\optslong{It is a routing protocol for sensor networks}{It is a link-layer protocol for RFID}{It is a request-response protocol that runs over UDP}{It is a lightweight publish-subscribe protocol that runs over TCP and uses a broker}
\end{question}

\begin{question}{H}
Match List I with List II.\\
List I (IoT protocol): A. MQTT  B. CoAP  C. 6LoWPAN  D. ZigBee\\
List II (Description): I. Adaptation layer for carrying IPv6 over IEEE 802.15.4  II. Low-power mesh networking standard built on IEEE 802.15.4  III. Publish-subscribe messaging over TCP  IV. RESTful request-response protocol over UDP
\optsshort{A-III, B-IV, C-I, D-II}{A-III, B-IV, C-II, D-I}{A-III, B-I, C-IV, D-II}{A-IV, B-III, C-I, D-II}
\end{question}

\begin{question}{M}
The default ports of MQTT and CoAP are respectively
\optsshort{TCP 1883 and UDP 5683}{TCP 80 and UDP 53}{TCP 8883 and TCP 443}{TCP 5683 and UDP 1883}
\end{question}

\begin{question}{M}
How many quality-of-service levels does MQTT define?
\optsshort{3}{2}{4}{5}
\end{question}

\begin{question}{H}
Given below are two statements: one is labelled as Assertion A and the other is labelled as Reason R.\\
Assertion A: IPv6 is considered essential for the large-scale deployment of IoT.\\
Reason R: The 128-bit address space of IPv6 can give a unique address to each of billions of devices.
\optslong{Both A and R are true and R is the correct explanation of A}{Both A and R are true but R is NOT the correct explanation of A}{A is true but R is false}{A is false but R is true}
\end{question}

\begin{question}{M}
Given below are two statements.\\
Statement I: Arduino Uno is a microcontroller-based board.\\
Statement II: Raspberry Pi is a single-board computer capable of running a full operating system such as Linux.
\optslong{Both Statement I and Statement II are true}{Both Statement I and Statement II are false}{Statement I is true but Statement II is false}{Statement I is false but Statement II is true}
\end{question}

\begin{question}{M}
Which of the following is a low-power wide-area network (LPWAN) technology for long-range IoT communication?
\optsshort{ZigBee}{NFC}{Bluetooth Low Energy}{LoRaWAN}
\end{question}

\begin{question}{M}
Which technology uses electromagnetic fields to automatically identify and track tags attached to objects?
\optsshort{OCR}{LiDAR}{RFID}{GPS}
\end{question}

\begin{question}{M}
A passive RFID tag
\optslong{is powered by the energy of the reader's radio signal}{has its own battery}{uses GPS to report location}{can be read from several kilometres}
\end{question}

\begin{question}{M}
A device that connects IoT devices using protocols such as ZigBee or BLE to the Internet, performing protocol translation, is called an
\optsshort{IoT gateway}{end node}{actuator}{RFID tag}
\end{question}

\begin{question}{M}
Direct communication between devices without human intervention, a precursor of IoT, is called
\optsshort{H2M interaction}{P2P file sharing}{B2B commerce}{M2M communication}
\end{question}

\begin{question}{M}
A virtual representation of a physical object or system that is kept synchronised with it through sensor data is called a
\optsshort{digital twin}{hologram}{data lake}{virtual machine}
\end{question}

\begin{question}{M}
Which of the following are major challenges of IoT?\\
A. Security and privacy  B. Interoperability  C. Energy constraints of devices  D. Scalability\\
Choose the correct answer from the options given below:
\optsshort{B, C and D only}{A, B, C and D}{A, B and C only}{A and B only}
\end{question}

\begin{question}{M}
Which of the following is NOT a characteristic of the Constrained Application Protocol (CoAP)?
\optslong{It is designed for constrained devices}{It requires a broker to deliver messages}{It runs over UDP}{It follows the REST model with methods such as GET and PUT}
\end{question}

\begin{question}{L}
In cloud computing, the term ``utility computing'' refers to
\optslong{free computing resources}{using utility software}{computing performed only by public utilities}{providing computing resources as a metered service, like electricity}
\end{question}

\begin{question}{M}
Which of the following is an example of Software as a Service?
\optslong{Salesforce CRM}{Heroku}{Google Compute Engine}{Microsoft Azure Virtual Machines}
\end{question}

\begin{question}{M}
Grid computing differs from cloud computing mainly in that a grid
\optslong{always belongs to one commercial provider}{is a federation of loosely coupled, geographically distributed and often heterogeneous resources applied to a common task}{cannot run parallel jobs}{offers on-demand self-service with pay-per-use billing}
\end{question}

\begin{question}{M}
In the cloud shared-responsibility model for IaaS, which of the following is the responsibility of the customer?
\optslong{Cooling and power}{Maintenance of the hypervisor}{Physical security of the data centre}{Patching the guest operating system and securing applications and data}
\end{question}

\begin{answers}
\ans{1}{4}{Gmail is ready-to-use software, App Engine a development platform, EC2 raw virtual machines.}
\ans{2}{3}{IaaS gives the consumer control from the operating system upwards.}
\ans{3}{3}{PaaS provides a managed platform; the consumer manages only applications and data.}
\ans{4}{3}{SaaS gives the least control and IaaS the most.}
\ans{5}{1}{These are the five essential characteristics in the NIST definition.}
\ans{6}{4}{The four NIST deployment models.}
\ans{7}{3}{Cloud bursting is a hybrid-cloud technique for handling peak loads.}
\ans{8}{1}{Type 1 is bare-metal (e.g. Xen, VMware ESXi, Hyper-V); Type 2 is hosted (e.g. VirtualBox, VMware Workstation).}
\ans{9}{4}{VirtualBox runs on top of a conventional operating system.}
\ans{10}{1}{Paravirtualization trades guest modification for better performance.}
\ans{11}{3}{It gives each VM the illusion of its own hardware.}
\ans{12}{1}{Sharing the kernel gives faster start-up and lower overhead.}
\ans{13}{1}{Live migration is used for load balancing and maintenance.}
\ans{14}{1}{Multi-tenancy is the basis of resource pooling in SaaS.}
\ans{15}{1}{Standard definitions of scale-up and scale-out.}
\ans{16}{4}{Scalability is the ability to grow; elasticity is growing and shrinking automatically.}
\ans{17}{2}{SLAs state metrics such as uptime percentage and response time.}
\ans{18}{3}{0.1\% of 30 $\times$ 24 $\times$ 60 = 43,200 minutes is 43.2 minutes.}
\ans{19}{4}{Objects in buckets, blocks as volumes, files in shared directories.}
\ans{20}{1}{Examples are Amazon RDS and Azure SQL Database.}
\ans{21}{1}{When a network partition occurs, a system must choose between consistency and availability.}
\ans{22}{4}{Consumers pay for measured usage instead of investing in infrastructure.}
\ans{23}{4}{Moving data and applications between providers can be difficult.}
\ans{24}{4}{AWS Lambda is a well-known example.}
\ans{25}{1}{It reduces latency and bandwidth, which matters for IoT applications.}
\ans{26}{3}{OpenStack provides compute (Nova), storage (Swift, Cinder) and networking (Neutron) services.}
\ans{27}{1}{Load balancing improves both utilisation and availability.}
\ans{28}{3}{Map tasks process key-value pairs and reduce tasks aggregate them.}
\ans{29}{1}{Consolidation is one of the main motivations for virtualization.}
\ans{30}{2}{Ashton used the term in 1999 in the context of RFID supply chains.}
\ans{31}{2}{Perception (sensing), network (transmission), application.}
\ans{32}{1}{Sensors measure; actuators act.}
\ans{33}{4}{Clients publish to topics and the broker delivers messages to subscribers.}
\ans{34}{1}{Standard roles of the four protocols.}
\ans{35}{1}{MQTT uses 1883 (8883 over TLS); CoAP uses UDP 5683.}
\ans{36}{1}{QoS 0 (at most once), QoS 1 (at least once), QoS 2 (exactly once).}
\ans{37}{1}{IPv4's 32-bit space is too small for the number of connected things.}
\ans{38}{1}{Arduino runs a single program on a microcontroller; Raspberry Pi has a microprocessor and an operating system.}
\ans{39}{4}{LoRaWAN and NB-IoT cover kilometres at low data rates; the others are short-range.}
\ans{40}{3}{A reader interrogates passive or active tags; RFID was an early enabler of IoT.}
\ans{41}{1}{Passive tags have no power source and a short read range.}
\ans{42}{1}{Gateways bridge constrained device networks and IP networks.}
\ans{43}{4}{Machine-to-machine communication.}
\ans{44}{1}{Digital twins are used for monitoring and simulation in industrial IoT.}
\ans{45}{2}{All four are well-recognised IoT challenges.}
\ans{46}{2}{CoAP is a direct client-server protocol; brokers are used by MQTT.}
\ans{47}{4}{Resources are packaged and billed according to usage.}
\ans{48}{1}{Salesforce delivers a finished application; Heroku is PaaS and the others are IaaS.}
\ans{49}{2}{Clouds are centrally managed, virtualized and offered as a metered service.}
\ans{50}{4}{The provider secures the cloud infrastructure; the customer secures what it runs in the cloud.}
\end{answers}
